% 完整 DMMR 模型图
% main.tex 可改为 % 完整 DMMR 模型图
% main.tex 可改为 % 完整 DMMR 模型图
% main.tex 可改为 % 完整 DMMR 模型图
% main.tex 可改为 \input{models/dmmr.tex} 来绘制该模型

\input{modules/eeg_signal.tex} % 引入EEG输入卡片
\input{modules/model_common.tex} % 引入通用节点和箭头样式
\input{modules/dmmr_components.tex} % 引入DMMR专用组件

\begin{tikzpicture}[
  node distance=8mm and 11mm
] % DMMR模型图开始
  % 全局间距参数：主要修改这些值即可调整布局
  \def\stepGap{4mm} % 预处理模块横向间距
  \def\preGap{4mm} % 预处理到预训练区域的间距
  \def\mmrGap{4mm} % Encoder到Decoder列、Decoder到重建列的间距
  \def\mixGap{8mm} % Mix到第二阶段Encoder的间距
  \def\mixEncGap{3mm} % 第二阶段Encoder到Decoder列的间距
  \def\mixDecGap{4mm} % 第二阶段多个Decoder的横向间距
  \def\mixReconGap{4mm} % 第二阶段Decoder到recon EEG的间距
  \def\mixBraceGap{1mm} % 第二阶段花括号与Decoder的间距
  \def\mixLabelGap{11mm} % Lrecon标签与花括号的间距
  \def\lossGap{4mm} % 重建输出到loss列的间距
  \def\decoderGap{8mm} % 多个Decoder的纵向间距
  \def\dgdannGap{10mm} % DG-DANN距离重建分支的垂直间距
  \def\dgdannInnerXSep{3mm} % DG-DANN外框横向内边距，pretrain会跟随自适应
  \def\dgdannInnerYSep{4mm} % DG-DANN外框纵向内边距，pretrain会跟随自适应
  \def\preShiftX{0mm} % 整个预训练区域的横向总偏移
  \def\preShiftY{0mm} % 整个预训练区域的纵向总偏移
  \def\fineGap{15mm} % Fine-tuning距离Pre-training的垂直间距
  \def\fineBaseX{-9mm} % Fine-tuning区域的独立横向锚点
  \def\fineBaseY{-40mm} % Fine-tuning区域的独立纵向锚点
  %
  % 输入和预处理流程：从左到右用相对定位排列
  \begin{scope}[shift={(\preShiftX,\preShiftY)}]
  \eegDataCard{rawSample}{Subject i}{dann source data} % 原始EEG样本
  \coordinate (raw-east) at (rawSample.east); % 原始样本连接点
  \node[dmmr prep, right=\stepGap of rawSample] (shuffleNode) {Time-step\\Shuffle}; % 时间步打乱
  \node[dmmr attention, right=\stepGap of shuffleNode] (abpNode) {ABP}; % Attention Based Pooling
  \node[dann source data, minimum width=24mm, minimum height=12mm, right=\stepGap of abpNode] (weightedSample) {}; % 加权样本外框
  \node[font=\scriptsize] at ([yshift=2.4mm]weightedSample.center) {Weighted Sample}; % 加权样本标签
  \begin{scope}[shift={([yshift=-2mm]weightedSample.center)}, scale=.62]
    \pic {eeg signal}; % 加权样本波形
  \end{scope}
  \coordinate (weighted-east) at (weightedSample.east); % 加权样本连接点
  \end{scope}
  %
  % DMMR Pre-training主体：后续模块相对weightedSample定位，因此会跟随preShift整体移动
  \node[dmmr encoder, right=\preGap of weightedSample] (encoderNode) {Encoder}; % 共享编码器
  \node[font=\scriptsize, text=black] (sharedFeatureText) at ([yshift=-4mm]encoderNode.south) {}; % 编码器说明
  \node[dmmr decoder, right=\mmrGap of encoderNode] (decI) {Decoder\\$D_{si}$}; % Subject i解码器
  \node[dmmr decoder, above=\decoderGap of decI] (decOne) {Decoder\\$D_{s1}$}; % Subject 1解码器
  \node[dmmr decoder, below=\decoderGap of decI] (decN) {Decoder\\$D_{sn}$}; % Subject n解码器
  \node[font=\scriptsize] at ($(decOne)!0.5!(decI)$) {$\vdots$}; % Decoder上方省略号
  \node[font=\scriptsize] at ($(decI)!0.5!(decN)$) {$\vdots$}; % Decoder下方省略号
  %
  \node[dmmr mix] (mixNode) at ([xshift=\mixGap]decI.east) {Mix}; % 第一阶段Decoder后接Mix节点
  \node[dmmr encoder, right=\mixEncGap of mixNode] (mixEncoder) {Encoder}; % 第二阶段编码器
  \node[dmmr decoder, right=\mixDecGap of mixEncoder] (mixDecI) {Decoder\\$D_{si}$}; % 第二阶段解码器
  \node[dmmr decoder, above=\decoderGap of mixDecI] (mixDecOne) {Decoder\\$D_{s1}$}; % 二阶段解码器1
  \node[dmmr decoder, below=\decoderGap of mixDecI] (mixDecN) {Decoder\\$D_{sn}$}; % 二阶段解码器n
  \node[font=\scriptsize] at ($(mixDecOne)!0.5!(mixDecI)$) {$\vdots$}; % 二阶段上方省略号
  \node[font=\scriptsize] at ($(mixDecI)!0.5!(mixDecN)$) {$\vdots$}; % 二阶段下方省略号
  \draw[dmmr recon edge] (decOne.east) -- (mixNode.west); % 第一阶段Decoder 1到Mix
  \draw[dmmr recon edge] (decI.east) -- (mixNode.west); % 第一阶段Decoder i到Mix
  \draw[dmmr recon edge] (decN.east) -- (mixNode.west); % 第一阶段Decoder n到Mix
  \draw[dmmr recon edge] (mixNode.east) -- (mixEncoder.west); % Mix到第二阶段Encoder
  \draw[dmmr recon edge] (mixEncoder.east) -- (mixDecOne.west); % 第二阶段Encoder到Decoder 1
  \draw[dmmr recon edge] (mixEncoder.east) -- (mixDecI.west); % 第二阶段Encoder到Decoder i
  \draw[dmmr recon edge] (mixEncoder.east) -- (mixDecN.west); % 第二阶段Encoder到Decoder n
  \begin{scope}[shift={([xshift=\mixReconGap]mixDecOne.east)}]
    \dmmrReconCard{reconOne}{\shortstack{recon $EEG_{1}$}}{dmmr recon output} % 第二阶段recon EEG 1
  \end{scope}
  \begin{scope}[shift={([xshift=\mixReconGap]mixDecI.east)}]
    \dmmrReconCard{reconI}{\shortstack{recon $EEG_{i}$}}{dmmr recon output} % 第二阶段recon EEG i
  \end{scope}
  \begin{scope}[shift={([xshift=\mixReconGap]mixDecN.east)}]
    \dmmrReconCard{reconN}{\shortstack{recon $EEG_{n}$}}{dmmr recon output} % 第二阶段recon EEG n
  \end{scope}
  \draw[dmmr recon edge] (mixDecOne.east) -- (reconOne.west); % 第二阶段Decoder 1到recon EEG 1
  \draw[dmmr recon edge] (mixDecI.east) -- (reconI.west); % 第二阶段Decoder i到recon EEG i
  \draw[dmmr recon edge] (mixDecN.east) -- (reconN.west); % 第二阶段Decoder n到recon EEG n
  \coordinate (reconBraceTop) at ([xshift=\mixBraceGap]reconOne.north east);
  \coordinate (reconBraceBottom) at ([xshift=\mixBraceGap]reconN.south east);
  \draw[
    decorate,
    decoration={brace, mirror, amplitude=3pt},
    draw=gray!70,
    line width=.5pt
  ] (reconBraceBottom) -- (reconBraceTop)
    node[midway, xshift=\mixLabelGap, dmmr loss] (reconLossLabel) {$L_{recon}$}; % 第二阶段recon EEG汇总loss方框
  \node[dmmr sum, right=\lossGap of reconLossLabel] (preSum) {$+$}; % Lrecon和Ladv相加
  \node[dmmr total loss, below=4mm of preSum] (preLoss) {$L_{pre}$}; % 总损失方框
  \coordinate (reconBraceMid) at ($(reconBraceTop)!0.5!(reconBraceBottom)$);
  \draw[dmmr recon edge] ([xshift=1.5mm]reconBraceMid) -- (reconLossLabel.west); % 从花括号右侧精确指向Lrecon的绿色箭头
  %
  % DG-DANN对抗分支：基于Decoder上方锚点整体上移
  \node[dann grl, minimum width=8mm, minimum height=6mm] (grlNode) at ([yshift=\dgdannGap]decOne.north) {GRL}; % 梯度反转层
  \node[dmmr subject discriminator, right=\stepGap of grlNode] (subjectDisc) {Subject\\Discriminator}; % 主体判别器
  \node[dmmr loss, right=\stepGap of subjectDisc] (advLoss) {$L_{adv}$}; % 对抗损失
  %
  % 主流程箭头
  \draw[dmmr edge] (raw-east) -- (shuffleNode.west); % 原始样本到时间步打乱
  \draw[dmmr edge] (shuffleNode.east) -- (abpNode.west); % 打乱样本到ABP
  \draw[dmmr edge] (abpNode.east) -- (weightedSample.west); % ABP到加权样本
  \draw[dmmr edge] (weighted-east) -- (encoderNode.west); % 加权样本到编码器
  %
  % 第一阶段编码器到Decoder分支
  \draw[dmmr recon edge] (encoderNode.east) -- (decOne.west); % 编码特征到Decoder 1
  \draw[dmmr recon edge] (encoderNode.east) -- (decI.west); % 编码特征到Decoder i
  \draw[dmmr recon edge] (encoderNode.east) -- (decN.west); % 编码特征到Decoder n
  \draw[dmmr recon edge] (reconLossLabel.east) -- (preSum.west); % Lrecon进入加法节点
  \draw[dmmr recon edge] (preSum.south) -- (preLoss.north); % 加法节点输出Lpre
  %
  % 对抗分支箭头
  \draw[dmmr adv edge] (encoderNode.north) |- (grlNode.west); % 编码特征进入GRL
  \draw[dmmr adv edge] (grlNode.east) -- (subjectDisc.west); % GRL到主体判别器
  \draw[dmmr adv edge] (subjectDisc.east) -- (advLoss.west); % 主体判别器到对抗损失
  \draw[dmmr adv edge] (advLoss.east) -| (preSum.north); % Ladv进入加法节点
  %
  % 背景分区：pretrainBox 的 fit 节点决定上方子图整体范围；fineInput 的锚点决定下方子图整体位置
  \begin{scope}[on background layer] % 背景层
    \node[
      inner xsep=\dgdannInnerXSep,
      inner ysep=\dgdannInnerYSep,
      fit=(grlNode)(subjectDisc)(advLoss)
    ] (dgdannFit) {}; % DG-DANN外框尺寸锚点，供Pre-training自适应包含
    \node[
      dmmr stage,
      inner xsep=3mm,
      inner ysep=3mm,
      fit=(rawSample)(shuffleNode)(abpNode)(weightedSample)(encoderNode)(sharedFeatureText)(decOne)(decI)(decN)(reconOne)(reconI)(reconN)(mixNode)(mixEncoder)(mixDecOne)(mixDecI)(mixDecN)(reconLossLabel)(preSum)(preLoss)(dgdannFit)
    ] (pretrainBox) {}; % 扩大后的Pre-training区域
    \node[
      dmmr stage,
      draw=gray!60,
      fill=gray!15,
      fill opacity=.90,
      inner xsep=\dgdannInnerXSep,
      inner ysep=\dgdannInnerYSep,
      fit=(grlNode)(subjectDisc)(advLoss)
    ] (dgdannBox) {}; % DG-DANN对抗区域
  \end{scope}
  \node[font=\scriptsize\bfseries, anchor=west] at ([xshift=1mm,yshift=-3mm]pretrainBox.north west) {Pre-training}; % 预训练标题
  \node[font=\fontsize{6pt}{7pt}\selectfont\bfseries, anchor=west] at ([xshift=1mm,yshift=-2mm]dgdannBox.north west) {DG-DANN}; % DG-DANN标题
  %
  % Fine-tuning分支：使用独立锚点，避免跟随pretrainBox一起移动
  \coordinate (fineOrigin) at (\fineBaseX, \fineBaseY); % Fine-tuning独立锚点
  \node[dmmr prep, anchor=west] (fineInput) at (fineOrigin) {EEG\\Sample}; % 微调输入
  \node[dmmr attention, right=\stepGap of fineInput] (fineABP) {ABP}; % 复用注意力层
  \node[dmmr encoder, right=\stepGap of fineABP] (fineEncoder) {Pretrained\\Encoder}; % 复用预训练编码器
  \node[dann head, right=\stepGap of fineEncoder] (fineClassifier) {Classifier\\64$\to$3}; % 微调分类头
  \node[dann output, right=\stepGap of fineClassifier] (finePred) {Predicted\\Label}; % 情绪预测
  \node[font=\scriptsize, text=black, anchor=west] at ([xshift=4mm]finePred.east) {}; % 微调说明
  %
  \draw[dmmr edge] (fineInput.east) -- (fineABP.west); % 微调输入到Attention
  \draw[dmmr edge] (fineABP.east) -- (fineEncoder.west); % Attention到Encoder
  \draw[dmmr edge] (fineEncoder.east) -- (fineClassifier.west); % Encoder到分类层
  \draw[dmmr edge] (fineClassifier.east) -- (finePred.west); % 分类层到情绪预测
  %
  \begin{scope}[on background layer] % 背景层
    \node[
      dmmr stage,
      fill=orange!5,
      inner xsep=3mm,
      inner ysep=3mm,
      fit=(fineInput)(fineABP)(fineEncoder)(fineClassifier)(finePred)
    ] (fineBox) {}; % 微调区域
  \end{scope}
  \node[font=\scriptsize\bfseries, anchor=west] at ([xshift=1mm,yshift=-3mm]fineBox.north west) {Fine-tuning}; % 微调标题
\end{tikzpicture}
 来绘制该模型

% 集中定义 EEG 信号波形模块
\providecommand{\eegDataCard}[3]{ % 节点名/标签/样式
  \node[#3] (#1) at (0, 0) {}; % 数据外框
  \node[font=\scriptsize] at ([yshift=2.4mm]#1.center) {#2}; % 数据标签
  \begin{scope}[shift={([yshift=-2mm]#1.center)}, scale=.62] % EEG波形位置
    \pic {eeg signal}; % EEG波形
  \end{scope}
}
%
\tikzset{
  % EEG白色背景样式
  eeg signal bg/.style={
    draw=gray!25, % 背景边框色
    fill=white, % 背景填充色
    rounded corners=1pt, % 背景轻微圆角
    line width=.25pt % 背景边框宽度
  },
  % EEG波形线样式
  eeg signal line/.style={
    draw=black!65,
    line width=.58pt,
    line cap=round,
    line join=round
  },
  % EEG信号波形
  pics/eeg signal/.style={
    code={ % 定义可复用EEG波形
      \path[eeg signal bg] (-.70, -.26) rectangle (.70, .30); % EEG白色背景
      \draw[eeg signal line]
        (-.62, .00) --
        (-.56, .01) --
        (-.50, -.02) --
        (-.45, .04) --
        (-.40, -.07) --
        (-.35, .05) --
        (-.30, -.02) --
        (-.25, -.15) --
        (-.21, .04) --
        (-.17, .18) --
        (-.13, .06) --
        (-.09, -.06) --
        (-.05, .02) --
        (.00, -.11) --
        (.05, -.04) --
        (.10, -.07) --
        (.16, .03) --
        (.21, -.03) --
        (.26, .08) --
        (.31, -.09) --
        (.36, .02) --
        (.42, -.04) --
        (.48, .03) --
        (.54, -.01) --
        (.58, .00) --
        (.62, .00); % EEG折线
    }
  }
}
 % 引入EEG输入卡片
% 集中定义模型图通用样式和通用组件
\tikzset{
  % 前向传播实线箭头
  dann edge/.style={
    -{Latex[length=1.5mm]}, % 箭头样式
    draw=teal!50!black, % 连线颜色
    line width=.55pt % 连线宽度
  },
  % 域适配/混淆损失虚线箭头
  dann confuse edge/.style={
    -{Latex[length=1.5mm]}, % 箭头样式
    draw=purple!65!black, % 连线颜色
    line width=.55pt, % 连线宽度
    % dashed % 虚线表示域混淆或分布对齐
  },
  % 输入数据节点基础样式
  dann data/.style={
    rounded corners=2pt, % 轻微圆角
    draw=blue!55!black, % 默认边框颜色
    fill=blue!9, % 默认填充颜色
    minimum width=18mm, % 最小宽度
    minimum height=10mm, % 最小高度
    align=center, % 居中换行
    font=\scriptsize % 使用较小字号
  },
  % Source数据节点样式
  dann source data/.style={
    dann data, % 继承数据节点基础样式
    draw=blue!75!black, % 源域边框色
    fill=blue!10 % 源域填充色
  },
  % Target数据节点样式
  dann target data/.style={
    dann data, % 继承数据节点基础样式
    draw=green!50!black, % 目标域边框色
    fill=green!12 % 目标域填充色
  },
  % 分类头节点样式
  dann head/.style={
    rounded corners=2pt, % 轻微圆角
    draw=orange!70!black, % 边框颜色
    fill=orange!13, % 填充颜色
    minimum width=18mm, % 最小宽度
    minimum height=9mm, % 最小高度
    align=center, % 居中换行
    font=\scriptsize % 使用较小字号
  },
  % GRL节点样式
  dann grl/.style={
    rounded corners=2pt, % 轻微圆角
    draw=purple!70!black, % 边框颜色
    fill=purple!12, % 填充颜色
    minimum width=16mm, % 最小宽度
    minimum height=9mm, % 最小高度
    align=center, % 居中换行
    font=\scriptsize % 使用较小加粗字号
  },
  % 域判别器节点样式
  dann domain/.style={
    rounded corners=2pt, % 轻微圆角
    draw=red!65!black, % 边框颜色
    fill=red!10, % 填充颜色
    minimum width=19mm, % 最小宽度
    minimum height=9mm, % 最小高度
    align=center, % 居中换行
    font=\scriptsize % 使用较小字号
  },
  % 输出节点样式
  dann output/.style={
    rounded corners=2pt, % 轻微圆角
    draw=gray!65, % 边框颜色
    fill=gray!8, % 填充颜色
    minimum width=17mm, % 最小宽度
    minimum height=8mm, % 最小高度
    align=center, % 居中换行
    font=\scriptsize % 使用较小字号
  },
  % 简化版共享CFE矩形样式
  dann cfe/.style={
    rounded corners=3pt, % 圆角
    draw=teal!60!black, % 边框颜色
    fill=teal!10, % 填充颜色
    minimum width=24mm, % 最小宽度
    minimum height=16mm, % 最小高度
    align=center, % 居中换行
    font=\scriptsize % 使用较小加粗字号
  },
  % Source EEG 输入组件
  pics/source eeg/.style={
    code={ % 定义Source EEG模块
      \eegDataCard{sourceData}{Source data}{dann source data} % 源域数据卡片
      \coordinate (source-east) at (sourceData.east); % 右侧连接点
    }
  },
  % Target EEG 输入组件
  pics/target eeg/.style={
    code={ % 定义Target EEG模块
      \eegDataCard{targetData}{Target data}{dann target data} % 目标域数据卡片
      \coordinate (target-east) at (targetData.east); % 右侧连接点
    }
  },
  % 简化版共享CFE组件
  pics/shared cfe/.style={
    code={ % 定义共享特征提取器
      \node[dann cfe] (sharedCFE) at (0, 0) {Shared CFE\\310$\to$64}; % 共享CFE矩形
      \node[font=\scriptsize, text=black] at (0, -.92) {Feature Extractor}; % CFE说明
      \coordinate (cfe-source-west) at ([yshift=2.5mm]sharedCFE.west); % 源域输入连接点
      \coordinate (cfe-target-west) at ([yshift=-2.5mm]sharedCFE.west); % 目标域输入连接点
      \coordinate (cfe-emo-east) at ([yshift=2.5mm]sharedCFE.east); % 情绪分支输出点
      \coordinate (cfe-domain-east) at ([yshift=-2.5mm]sharedCFE.east); % 域分支输出点
    }
  },
  % 情绪分类器组件
  pics/emotion classifier/.style={
    code={ % 定义情绪分类器
      \node[dann head] (emoHead) at (0, 0) {Emotion\\Classifier}; % 线性分类层
      \node[dann output] (emoOut) at (2.25, 0) {Predicted\\Label}; % 情感预测
      \draw[dann edge] (emoHead) -- (emoOut); % 分类器到输出
      \coordinate (emo-west) at (emoHead.west); % 左侧连接点
    }
  },
  % GRL组件
  pics/grl/.style={
    code={ % 定义梯度反转层
      \node[dann grl] (grlNode) at (0, 0) {GRL}; % 梯度反转层
      \coordinate (grl-west) at (grlNode.west); % 左侧连接点
      \coordinate (grl-east) at (grlNode.east); % 右侧连接点
    }
  },
  % 域判别器组件
  pics/domain discriminator/.style={
    code={ % 定义域判别器
      \node[dann domain] (domainHead) at (0, 0) {Domain\\Discriminator}; % 域判别器
      \node[dann output] (domainOut) at (2.45, 0) {Source / Target\\Confusion}; % 域混淆输出
      \draw[dann confuse edge] (domainHead) -- (domainOut); % 域判别到混淆
      \coordinate (domain-west) at (domainHead.west); % 左侧连接点
    }
  }
}

% 通用源域/目标域特征图例
\providecommand{\domainFeatureLegend}[2]{ % 锚点节点/下方距离
  \matrix (legendContent) [
    below=#2 of #1,
    column sep=2mm,
    inner sep=0pt,
    nodes={inner sep=0pt, anchor=center},
    ampersand replacement=\&
  ] {
    \node[minimum size=6mm] (legendSourceIcon) {}; \&
    \node[anchor=west, font=\scriptsize, text=gray!95] (legendSourceText) {Source feature}; \&
    \node[minimum width=5mm] {}; \&
    \node[minimum size=6mm] (legendTargetIcon) {}; \&
    \node[anchor=west, font=\scriptsize, text=gray!95] (legendTargetText) {Target feature}; \\
  }; % 图例内容矩阵
  \begin{scope}[shift={(legendSourceIcon.center)}, scale=.30, transform shape, name prefix=legendSource-]
    \pic {source feature cloud}; % 源域特征图例
  \end{scope}
  \begin{scope}[shift={(legendTargetIcon.center)}, scale=.30, transform shape, name prefix=legendTarget-]
    \pic {target feature cloud}; % 目标域特征图例
  \end{scope}
  \begin{scope}[on background layer]
    \node[
      rounded corners=4pt,
      draw=gray!85,
      dashed,
      inner xsep=5mm,
      inner ysep=2mm,
      fit=(legendContent)
    ] (legendBox) {}; % 图例外框
  \end{scope}
}
 % 引入通用节点和箭头样式
% 集中定义 DMMR 专用组件
\providecommand{\dmmrReconCard}[3]{% 节点名/标签/样式
  \node[#3, anchor=west, minimum width=16mm, minimum height=10mm, inner sep=0pt] (#1) at (0, 0) {}; % 重建输出外框
  \node[font=\scriptsize, align=center, anchor=north] at ([yshift=-1mm]#1.north) {#2}; % 重建输出标签
  \begin{scope}
    \clip ([xshift=1mm,yshift=1mm]#1.south west) rectangle ([xshift=-1mm,yshift=-3mm]#1.north east); % 限制波形始终在框内
    \begin{scope}[shift={([yshift=2.6mm]#1.south)}, scale=.38] % 重建波形位置
      \pic {eeg signal}; % EEG波形
    \end{scope}
  \end{scope}
  \coordinate (#1-west) at (#1.west); % 左侧连接点
  \coordinate (#1-east) at (#1.east); % 右侧连接点
}

\tikzset{
  % DMMR预处理节点样式
  dmmr prep/.style={
    rounded corners=2pt, % 轻微圆角
    draw=gray!70, % 边框颜色
    fill=gray!8, % 填充颜色
    minimum width=18mm, % 节点宽度
    minimum height=9mm, % 节点高度
    align=center, % 文本居中
    font=\scriptsize % 使用较小字号
  },
  % DMMR注意力节点样式
  dmmr attention/.style={
    rounded corners=2pt, % 轻微圆角
    draw=purple!55!black, % 边框颜色
    fill=purple!12, % 填充颜色
    minimum width=8mm, % 节点宽度
    minimum height=18mm, % 节点高度
    align=center, % 文本居中
    font=\scriptsize % 使用较小字号
  },
  % DMMR编码器节点样式
  dmmr encoder/.style={
    rounded corners=7pt, % 胶囊形圆角
    draw=teal!60!black, % 边框颜色
    fill=cyan!18, % 填充颜色
    minimum width=10mm, % 节点宽度
    minimum height=6mm, % 节点高度
    align=center, % 文本居中
    font=\scriptsize % 使用较小字号
  },
  % DMMR解码器节点样式
  dmmr decoder/.style={
    rounded corners=7pt, % 胶囊形圆角
    draw=green!45!black, % 边框颜色
    fill=green!14, % 填充颜色
    minimum width=12mm, % 节点宽度
    minimum height=8mm, % 节点高度
    align=center, % 文本居中
    font=\scriptsize % 使用较小字号
  },
  % DMMR域判别器节点样式
  dmmr subject discriminator/.style={
    rounded corners=7pt, % 胶囊形圆角
    draw=blue!55!black, % 边框颜色
    fill=blue!18, % 填充颜色
    minimum width=20mm, % 节点宽度
    minimum height=8mm, % 节点高度
    align=center, % 文本居中
    font=\scriptsize % 使用较小字号
  },
  % DMMR阶段背景样式
  dmmr stage/.style={
    rounded corners=4pt, % 背景圆角
    draw=gray!70, % 背景边框
    dashed, % 虚线边框
    fill=green!7, % 默认浅绿背景
    inner sep=3mm % 背景内边距
  },
  % DMMR损失节点样式
  dmmr loss/.style={
    rounded corners=2pt, % 轻微圆角
    draw=orange!75!black, % loss边框
    fill=orange!6, % loss背景
    minimum width=8mm, % 节点宽度
    minimum height=6mm, % 节点高度
    align=center, % 文本居中
    font=\scriptsize % 使用较小字号
  },
  % DMMR总损失样式
  dmmr total loss/.style={
    rounded corners=2pt, % 轻微圆角
    draw=red!70!black, % 总损失边框
    fill=red!10, % 总损失背景
    minimum width=10mm, % 节点宽度
    minimum height=6mm, % 节点高度
    align=center, % 文本居中
    font=\scriptsize % 使用较小字号
  },
  % DMMR混合节点样式
  dmmr mix/.style={
    circle,
    draw=gray!70,
    fill=white,
    minimum size=6mm,
    inner sep=0pt,
    font=\scriptsize
  },
  % DMMR损失相加节点样式
  dmmr sum/.style={
    circle, % 圆形节点
    draw=gray!70, % 边框颜色
    fill=white, % 背景颜色
    minimum size=4mm, % 圆圈大小
    inner sep=0pt, % 无内边距
    font=\scriptsize\bfseries % 加号字号
  },
  % DMMR重建输出样式
  dmmr recon output/.style={
    rounded corners=1.5pt, % 轻微圆角
    draw=gray!50, % 边框颜色
    fill=white, % 背景颜色
    minimum width=16mm, % 节点宽度
    minimum height=11mm, % 节点高度
    inner xsep=1mm, % 横向内边距
    inner ysep=1mm, % 纵向内边距
    align=center, % 文本居中
    font=\scriptsize % 使用较小字号
  },
  % DMMR前向箭头
  dmmr edge/.style={
    -{Latex[length=1.5mm]}, % 箭头样式
    draw=teal!50!black, % 箭头颜色
    line width=.55pt % 箭头线宽
  },
  % DMMR对抗箭头
  dmmr adv edge/.style={
    -{Latex[length=1.5mm]}, % 箭头样式
    draw=purple!65!black, % 对抗分支颜色
    line width=.55pt % 箭头线宽
  },
  % DMMR重建箭头
  dmmr recon edge/.style={
    -{Latex[length=1.5mm]}, % 箭头样式
    draw=green!45!black, % 重建分支颜色
    line width=.55pt % 箭头线宽
  },
  % 小型EEG输出样式
  dmmr mini eeg box/.style={
    rounded corners=1pt, % 轻微圆角
    draw=gray!55, % 边框颜色
    fill=white, % 背景颜色
    minimum width=12mm, % 节点宽度
    minimum height=7mm, % 节点高度
    inner sep=0pt % 无内边距
  },
  % 小型EEG输出组件
  pics/dmmr mini eeg/.style={
    code={ % 定义重建EEG小图
      \node[dmmr mini eeg box] (miniEEGBox) at (0, 0) {}; % 小图外框
      \begin{scope}[scale=.43, transform shape]
        \pic at (0, 0) {eeg signal}; % 复用EEG波形
      \end{scope}
      \coordinate (miniEEG-west) at (miniEEGBox.west); % 左侧连接点
      \coordinate (miniEEG-east) at (miniEEGBox.east); % 右侧连接点
    }
  },
  % DMMR预训练总损失组件
  pics/dmmr pretrain loss/.style={
    code={ % 定义预训练总损失
      \node[dmmr loss] (reconLoss) at (0, .45) {$L_{recon}$}; % 重建损失
      \node[dmmr loss] (advLoss) at (0, -.45) {$L_{adv}$}; % 对抗损失
      \node[
        font=\scriptsize,
        align=center
      ] (preLossText) at (0, -1.20) {$L_{pre}=L_{recon}+L_{adv}$}; % 总预训练损失
      \coordinate (preLoss-recon-west) at (reconLoss.west); % 重建损失入口
      \coordinate (preLoss-adv-west) at (advLoss.west); % 对抗损失入口
    }
  }
}
 % 引入DMMR专用组件

\begin{tikzpicture}[
  node distance=8mm and 11mm
] % DMMR模型图开始
  % 全局间距参数：主要修改这些值即可调整布局
  \def\stepGap{4mm} % 预处理模块横向间距
  \def\preGap{4mm} % 预处理到预训练区域的间距
  \def\mmrGap{4mm} % Encoder到Decoder列、Decoder到重建列的间距
  \def\mixGap{8mm} % Mix到第二阶段Encoder的间距
  \def\mixEncGap{3mm} % 第二阶段Encoder到Decoder列的间距
  \def\mixDecGap{4mm} % 第二阶段多个Decoder的横向间距
  \def\mixReconGap{4mm} % 第二阶段Decoder到recon EEG的间距
  \def\mixBraceGap{1mm} % 第二阶段花括号与Decoder的间距
  \def\mixLabelGap{11mm} % Lrecon标签与花括号的间距
  \def\lossGap{4mm} % 重建输出到loss列的间距
  \def\decoderGap{8mm} % 多个Decoder的纵向间距
  \def\dgdannGap{10mm} % DG-DANN距离重建分支的垂直间距
  \def\dgdannInnerXSep{3mm} % DG-DANN外框横向内边距，pretrain会跟随自适应
  \def\dgdannInnerYSep{4mm} % DG-DANN外框纵向内边距，pretrain会跟随自适应
  \def\preShiftX{0mm} % 整个预训练区域的横向总偏移
  \def\preShiftY{0mm} % 整个预训练区域的纵向总偏移
  \def\fineGap{15mm} % Fine-tuning距离Pre-training的垂直间距
  \def\fineBaseX{-9mm} % Fine-tuning区域的独立横向锚点
  \def\fineBaseY{-40mm} % Fine-tuning区域的独立纵向锚点
  %
  % 输入和预处理流程：从左到右用相对定位排列
  \begin{scope}[shift={(\preShiftX,\preShiftY)}]
  \eegDataCard{rawSample}{Subject i}{dann source data} % 原始EEG样本
  \coordinate (raw-east) at (rawSample.east); % 原始样本连接点
  \node[dmmr prep, right=\stepGap of rawSample] (shuffleNode) {Time-step\\Shuffle}; % 时间步打乱
  \node[dmmr attention, right=\stepGap of shuffleNode] (abpNode) {ABP}; % Attention Based Pooling
  \node[dann source data, minimum width=24mm, minimum height=12mm, right=\stepGap of abpNode] (weightedSample) {}; % 加权样本外框
  \node[font=\scriptsize] at ([yshift=2.4mm]weightedSample.center) {Weighted Sample}; % 加权样本标签
  \begin{scope}[shift={([yshift=-2mm]weightedSample.center)}, scale=.62]
    \pic {eeg signal}; % 加权样本波形
  \end{scope}
  \coordinate (weighted-east) at (weightedSample.east); % 加权样本连接点
  \end{scope}
  %
  % DMMR Pre-training主体：后续模块相对weightedSample定位，因此会跟随preShift整体移动
  \node[dmmr encoder, right=\preGap of weightedSample] (encoderNode) {Encoder}; % 共享编码器
  \node[font=\scriptsize, text=black] (sharedFeatureText) at ([yshift=-4mm]encoderNode.south) {}; % 编码器说明
  \node[dmmr decoder, right=\mmrGap of encoderNode] (decI) {Decoder\\$D_{si}$}; % Subject i解码器
  \node[dmmr decoder, above=\decoderGap of decI] (decOne) {Decoder\\$D_{s1}$}; % Subject 1解码器
  \node[dmmr decoder, below=\decoderGap of decI] (decN) {Decoder\\$D_{sn}$}; % Subject n解码器
  \node[font=\scriptsize] at ($(decOne)!0.5!(decI)$) {$\vdots$}; % Decoder上方省略号
  \node[font=\scriptsize] at ($(decI)!0.5!(decN)$) {$\vdots$}; % Decoder下方省略号
  %
  \node[dmmr mix] (mixNode) at ([xshift=\mixGap]decI.east) {Mix}; % 第一阶段Decoder后接Mix节点
  \node[dmmr encoder, right=\mixEncGap of mixNode] (mixEncoder) {Encoder}; % 第二阶段编码器
  \node[dmmr decoder, right=\mixDecGap of mixEncoder] (mixDecI) {Decoder\\$D_{si}$}; % 第二阶段解码器
  \node[dmmr decoder, above=\decoderGap of mixDecI] (mixDecOne) {Decoder\\$D_{s1}$}; % 二阶段解码器1
  \node[dmmr decoder, below=\decoderGap of mixDecI] (mixDecN) {Decoder\\$D_{sn}$}; % 二阶段解码器n
  \node[font=\scriptsize] at ($(mixDecOne)!0.5!(mixDecI)$) {$\vdots$}; % 二阶段上方省略号
  \node[font=\scriptsize] at ($(mixDecI)!0.5!(mixDecN)$) {$\vdots$}; % 二阶段下方省略号
  \draw[dmmr recon edge] (decOne.east) -- (mixNode.west); % 第一阶段Decoder 1到Mix
  \draw[dmmr recon edge] (decI.east) -- (mixNode.west); % 第一阶段Decoder i到Mix
  \draw[dmmr recon edge] (decN.east) -- (mixNode.west); % 第一阶段Decoder n到Mix
  \draw[dmmr recon edge] (mixNode.east) -- (mixEncoder.west); % Mix到第二阶段Encoder
  \draw[dmmr recon edge] (mixEncoder.east) -- (mixDecOne.west); % 第二阶段Encoder到Decoder 1
  \draw[dmmr recon edge] (mixEncoder.east) -- (mixDecI.west); % 第二阶段Encoder到Decoder i
  \draw[dmmr recon edge] (mixEncoder.east) -- (mixDecN.west); % 第二阶段Encoder到Decoder n
  \begin{scope}[shift={([xshift=\mixReconGap]mixDecOne.east)}]
    \dmmrReconCard{reconOne}{\shortstack{recon $EEG_{1}$}}{dmmr recon output} % 第二阶段recon EEG 1
  \end{scope}
  \begin{scope}[shift={([xshift=\mixReconGap]mixDecI.east)}]
    \dmmrReconCard{reconI}{\shortstack{recon $EEG_{i}$}}{dmmr recon output} % 第二阶段recon EEG i
  \end{scope}
  \begin{scope}[shift={([xshift=\mixReconGap]mixDecN.east)}]
    \dmmrReconCard{reconN}{\shortstack{recon $EEG_{n}$}}{dmmr recon output} % 第二阶段recon EEG n
  \end{scope}
  \draw[dmmr recon edge] (mixDecOne.east) -- (reconOne.west); % 第二阶段Decoder 1到recon EEG 1
  \draw[dmmr recon edge] (mixDecI.east) -- (reconI.west); % 第二阶段Decoder i到recon EEG i
  \draw[dmmr recon edge] (mixDecN.east) -- (reconN.west); % 第二阶段Decoder n到recon EEG n
  \coordinate (reconBraceTop) at ([xshift=\mixBraceGap]reconOne.north east);
  \coordinate (reconBraceBottom) at ([xshift=\mixBraceGap]reconN.south east);
  \draw[
    decorate,
    decoration={brace, mirror, amplitude=3pt},
    draw=gray!70,
    line width=.5pt
  ] (reconBraceBottom) -- (reconBraceTop)
    node[midway, xshift=\mixLabelGap, dmmr loss] (reconLossLabel) {$L_{recon}$}; % 第二阶段recon EEG汇总loss方框
  \node[dmmr sum, right=\lossGap of reconLossLabel] (preSum) {$+$}; % Lrecon和Ladv相加
  \node[dmmr total loss, below=4mm of preSum] (preLoss) {$L_{pre}$}; % 总损失方框
  \coordinate (reconBraceMid) at ($(reconBraceTop)!0.5!(reconBraceBottom)$);
  \draw[dmmr recon edge] ([xshift=1.5mm]reconBraceMid) -- (reconLossLabel.west); % 从花括号右侧精确指向Lrecon的绿色箭头
  %
  % DG-DANN对抗分支：基于Decoder上方锚点整体上移
  \node[dann grl, minimum width=8mm, minimum height=6mm] (grlNode) at ([yshift=\dgdannGap]decOne.north) {GRL}; % 梯度反转层
  \node[dmmr subject discriminator, right=\stepGap of grlNode] (subjectDisc) {Subject\\Discriminator}; % 主体判别器
  \node[dmmr loss, right=\stepGap of subjectDisc] (advLoss) {$L_{adv}$}; % 对抗损失
  %
  % 主流程箭头
  \draw[dmmr edge] (raw-east) -- (shuffleNode.west); % 原始样本到时间步打乱
  \draw[dmmr edge] (shuffleNode.east) -- (abpNode.west); % 打乱样本到ABP
  \draw[dmmr edge] (abpNode.east) -- (weightedSample.west); % ABP到加权样本
  \draw[dmmr edge] (weighted-east) -- (encoderNode.west); % 加权样本到编码器
  %
  % 第一阶段编码器到Decoder分支
  \draw[dmmr recon edge] (encoderNode.east) -- (decOne.west); % 编码特征到Decoder 1
  \draw[dmmr recon edge] (encoderNode.east) -- (decI.west); % 编码特征到Decoder i
  \draw[dmmr recon edge] (encoderNode.east) -- (decN.west); % 编码特征到Decoder n
  \draw[dmmr recon edge] (reconLossLabel.east) -- (preSum.west); % Lrecon进入加法节点
  \draw[dmmr recon edge] (preSum.south) -- (preLoss.north); % 加法节点输出Lpre
  %
  % 对抗分支箭头
  \draw[dmmr adv edge] (encoderNode.north) |- (grlNode.west); % 编码特征进入GRL
  \draw[dmmr adv edge] (grlNode.east) -- (subjectDisc.west); % GRL到主体判别器
  \draw[dmmr adv edge] (subjectDisc.east) -- (advLoss.west); % 主体判别器到对抗损失
  \draw[dmmr adv edge] (advLoss.east) -| (preSum.north); % Ladv进入加法节点
  %
  % 背景分区：pretrainBox 的 fit 节点决定上方子图整体范围；fineInput 的锚点决定下方子图整体位置
  \begin{scope}[on background layer] % 背景层
    \node[
      inner xsep=\dgdannInnerXSep,
      inner ysep=\dgdannInnerYSep,
      fit=(grlNode)(subjectDisc)(advLoss)
    ] (dgdannFit) {}; % DG-DANN外框尺寸锚点，供Pre-training自适应包含
    \node[
      dmmr stage,
      inner xsep=3mm,
      inner ysep=3mm,
      fit=(rawSample)(shuffleNode)(abpNode)(weightedSample)(encoderNode)(sharedFeatureText)(decOne)(decI)(decN)(reconOne)(reconI)(reconN)(mixNode)(mixEncoder)(mixDecOne)(mixDecI)(mixDecN)(reconLossLabel)(preSum)(preLoss)(dgdannFit)
    ] (pretrainBox) {}; % 扩大后的Pre-training区域
    \node[
      dmmr stage,
      draw=gray!60,
      fill=gray!15,
      fill opacity=.90,
      inner xsep=\dgdannInnerXSep,
      inner ysep=\dgdannInnerYSep,
      fit=(grlNode)(subjectDisc)(advLoss)
    ] (dgdannBox) {}; % DG-DANN对抗区域
  \end{scope}
  \node[font=\scriptsize\bfseries, anchor=west] at ([xshift=1mm,yshift=-3mm]pretrainBox.north west) {Pre-training}; % 预训练标题
  \node[font=\fontsize{6pt}{7pt}\selectfont\bfseries, anchor=west] at ([xshift=1mm,yshift=-2mm]dgdannBox.north west) {DG-DANN}; % DG-DANN标题
  %
  % Fine-tuning分支：使用独立锚点，避免跟随pretrainBox一起移动
  \coordinate (fineOrigin) at (\fineBaseX, \fineBaseY); % Fine-tuning独立锚点
  \node[dmmr prep, anchor=west] (fineInput) at (fineOrigin) {EEG\\Sample}; % 微调输入
  \node[dmmr attention, right=\stepGap of fineInput] (fineABP) {ABP}; % 复用注意力层
  \node[dmmr encoder, right=\stepGap of fineABP] (fineEncoder) {Pretrained\\Encoder}; % 复用预训练编码器
  \node[dann head, right=\stepGap of fineEncoder] (fineClassifier) {Classifier\\64$\to$3}; % 微调分类头
  \node[dann output, right=\stepGap of fineClassifier] (finePred) {Predicted\\Label}; % 情绪预测
  \node[font=\scriptsize, text=black, anchor=west] at ([xshift=4mm]finePred.east) {}; % 微调说明
  %
  \draw[dmmr edge] (fineInput.east) -- (fineABP.west); % 微调输入到Attention
  \draw[dmmr edge] (fineABP.east) -- (fineEncoder.west); % Attention到Encoder
  \draw[dmmr edge] (fineEncoder.east) -- (fineClassifier.west); % Encoder到分类层
  \draw[dmmr edge] (fineClassifier.east) -- (finePred.west); % 分类层到情绪预测
  %
  \begin{scope}[on background layer] % 背景层
    \node[
      dmmr stage,
      fill=orange!5,
      inner xsep=3mm,
      inner ysep=3mm,
      fit=(fineInput)(fineABP)(fineEncoder)(fineClassifier)(finePred)
    ] (fineBox) {}; % 微调区域
  \end{scope}
  \node[font=\scriptsize\bfseries, anchor=west] at ([xshift=1mm,yshift=-3mm]fineBox.north west) {Fine-tuning}; % 微调标题
\end{tikzpicture}
 来绘制该模型

% 集中定义 EEG 信号波形模块
\providecommand{\eegDataCard}[3]{ % 节点名/标签/样式
  \node[#3] (#1) at (0, 0) {}; % 数据外框
  \node[font=\scriptsize] at ([yshift=2.4mm]#1.center) {#2}; % 数据标签
  \begin{scope}[shift={([yshift=-2mm]#1.center)}, scale=.62] % EEG波形位置
    \pic {eeg signal}; % EEG波形
  \end{scope}
}
%
\tikzset{
  % EEG白色背景样式
  eeg signal bg/.style={
    draw=gray!25, % 背景边框色
    fill=white, % 背景填充色
    rounded corners=1pt, % 背景轻微圆角
    line width=.25pt % 背景边框宽度
  },
  % EEG波形线样式
  eeg signal line/.style={
    draw=black!65,
    line width=.58pt,
    line cap=round,
    line join=round
  },
  % EEG信号波形
  pics/eeg signal/.style={
    code={ % 定义可复用EEG波形
      \path[eeg signal bg] (-.70, -.26) rectangle (.70, .30); % EEG白色背景
      \draw[eeg signal line]
        (-.62, .00) --
        (-.56, .01) --
        (-.50, -.02) --
        (-.45, .04) --
        (-.40, -.07) --
        (-.35, .05) --
        (-.30, -.02) --
        (-.25, -.15) --
        (-.21, .04) --
        (-.17, .18) --
        (-.13, .06) --
        (-.09, -.06) --
        (-.05, .02) --
        (.00, -.11) --
        (.05, -.04) --
        (.10, -.07) --
        (.16, .03) --
        (.21, -.03) --
        (.26, .08) --
        (.31, -.09) --
        (.36, .02) --
        (.42, -.04) --
        (.48, .03) --
        (.54, -.01) --
        (.58, .00) --
        (.62, .00); % EEG折线
    }
  }
}
 % 引入EEG输入卡片
% 集中定义模型图通用样式和通用组件
\tikzset{
  % 前向传播实线箭头
  dann edge/.style={
    -{Latex[length=1.5mm]}, % 箭头样式
    draw=teal!50!black, % 连线颜色
    line width=.55pt % 连线宽度
  },
  % 域适配/混淆损失虚线箭头
  dann confuse edge/.style={
    -{Latex[length=1.5mm]}, % 箭头样式
    draw=purple!65!black, % 连线颜色
    line width=.55pt, % 连线宽度
    % dashed % 虚线表示域混淆或分布对齐
  },
  % 输入数据节点基础样式
  dann data/.style={
    rounded corners=2pt, % 轻微圆角
    draw=blue!55!black, % 默认边框颜色
    fill=blue!9, % 默认填充颜色
    minimum width=18mm, % 最小宽度
    minimum height=10mm, % 最小高度
    align=center, % 居中换行
    font=\scriptsize % 使用较小字号
  },
  % Source数据节点样式
  dann source data/.style={
    dann data, % 继承数据节点基础样式
    draw=blue!75!black, % 源域边框色
    fill=blue!10 % 源域填充色
  },
  % Target数据节点样式
  dann target data/.style={
    dann data, % 继承数据节点基础样式
    draw=green!50!black, % 目标域边框色
    fill=green!12 % 目标域填充色
  },
  % 分类头节点样式
  dann head/.style={
    rounded corners=2pt, % 轻微圆角
    draw=orange!70!black, % 边框颜色
    fill=orange!13, % 填充颜色
    minimum width=18mm, % 最小宽度
    minimum height=9mm, % 最小高度
    align=center, % 居中换行
    font=\scriptsize % 使用较小字号
  },
  % GRL节点样式
  dann grl/.style={
    rounded corners=2pt, % 轻微圆角
    draw=purple!70!black, % 边框颜色
    fill=purple!12, % 填充颜色
    minimum width=16mm, % 最小宽度
    minimum height=9mm, % 最小高度
    align=center, % 居中换行
    font=\scriptsize % 使用较小加粗字号
  },
  % 域判别器节点样式
  dann domain/.style={
    rounded corners=2pt, % 轻微圆角
    draw=red!65!black, % 边框颜色
    fill=red!10, % 填充颜色
    minimum width=19mm, % 最小宽度
    minimum height=9mm, % 最小高度
    align=center, % 居中换行
    font=\scriptsize % 使用较小字号
  },
  % 输出节点样式
  dann output/.style={
    rounded corners=2pt, % 轻微圆角
    draw=gray!65, % 边框颜色
    fill=gray!8, % 填充颜色
    minimum width=17mm, % 最小宽度
    minimum height=8mm, % 最小高度
    align=center, % 居中换行
    font=\scriptsize % 使用较小字号
  },
  % 简化版共享CFE矩形样式
  dann cfe/.style={
    rounded corners=3pt, % 圆角
    draw=teal!60!black, % 边框颜色
    fill=teal!10, % 填充颜色
    minimum width=24mm, % 最小宽度
    minimum height=16mm, % 最小高度
    align=center, % 居中换行
    font=\scriptsize % 使用较小加粗字号
  },
  % Source EEG 输入组件
  pics/source eeg/.style={
    code={ % 定义Source EEG模块
      \eegDataCard{sourceData}{Source data}{dann source data} % 源域数据卡片
      \coordinate (source-east) at (sourceData.east); % 右侧连接点
    }
  },
  % Target EEG 输入组件
  pics/target eeg/.style={
    code={ % 定义Target EEG模块
      \eegDataCard{targetData}{Target data}{dann target data} % 目标域数据卡片
      \coordinate (target-east) at (targetData.east); % 右侧连接点
    }
  },
  % 简化版共享CFE组件
  pics/shared cfe/.style={
    code={ % 定义共享特征提取器
      \node[dann cfe] (sharedCFE) at (0, 0) {Shared CFE\\310$\to$64}; % 共享CFE矩形
      \node[font=\scriptsize, text=black] at (0, -.92) {Feature Extractor}; % CFE说明
      \coordinate (cfe-source-west) at ([yshift=2.5mm]sharedCFE.west); % 源域输入连接点
      \coordinate (cfe-target-west) at ([yshift=-2.5mm]sharedCFE.west); % 目标域输入连接点
      \coordinate (cfe-emo-east) at ([yshift=2.5mm]sharedCFE.east); % 情绪分支输出点
      \coordinate (cfe-domain-east) at ([yshift=-2.5mm]sharedCFE.east); % 域分支输出点
    }
  },
  % 情绪分类器组件
  pics/emotion classifier/.style={
    code={ % 定义情绪分类器
      \node[dann head] (emoHead) at (0, 0) {Emotion\\Classifier}; % 线性分类层
      \node[dann output] (emoOut) at (2.25, 0) {Predicted\\Label}; % 情感预测
      \draw[dann edge] (emoHead) -- (emoOut); % 分类器到输出
      \coordinate (emo-west) at (emoHead.west); % 左侧连接点
    }
  },
  % GRL组件
  pics/grl/.style={
    code={ % 定义梯度反转层
      \node[dann grl] (grlNode) at (0, 0) {GRL}; % 梯度反转层
      \coordinate (grl-west) at (grlNode.west); % 左侧连接点
      \coordinate (grl-east) at (grlNode.east); % 右侧连接点
    }
  },
  % 域判别器组件
  pics/domain discriminator/.style={
    code={ % 定义域判别器
      \node[dann domain] (domainHead) at (0, 0) {Domain\\Discriminator}; % 域判别器
      \node[dann output] (domainOut) at (2.45, 0) {Source / Target\\Confusion}; % 域混淆输出
      \draw[dann confuse edge] (domainHead) -- (domainOut); % 域判别到混淆
      \coordinate (domain-west) at (domainHead.west); % 左侧连接点
    }
  }
}

% 通用源域/目标域特征图例
\providecommand{\domainFeatureLegend}[2]{ % 锚点节点/下方距离
  \matrix (legendContent) [
    below=#2 of #1,
    column sep=2mm,
    inner sep=0pt,
    nodes={inner sep=0pt, anchor=center},
    ampersand replacement=\&
  ] {
    \node[minimum size=6mm] (legendSourceIcon) {}; \&
    \node[anchor=west, font=\scriptsize, text=gray!95] (legendSourceText) {Source feature}; \&
    \node[minimum width=5mm] {}; \&
    \node[minimum size=6mm] (legendTargetIcon) {}; \&
    \node[anchor=west, font=\scriptsize, text=gray!95] (legendTargetText) {Target feature}; \\
  }; % 图例内容矩阵
  \begin{scope}[shift={(legendSourceIcon.center)}, scale=.30, transform shape, name prefix=legendSource-]
    \pic {source feature cloud}; % 源域特征图例
  \end{scope}
  \begin{scope}[shift={(legendTargetIcon.center)}, scale=.30, transform shape, name prefix=legendTarget-]
    \pic {target feature cloud}; % 目标域特征图例
  \end{scope}
  \begin{scope}[on background layer]
    \node[
      rounded corners=4pt,
      draw=gray!85,
      dashed,
      inner xsep=5mm,
      inner ysep=2mm,
      fit=(legendContent)
    ] (legendBox) {}; % 图例外框
  \end{scope}
}
 % 引入通用节点和箭头样式
% 集中定义 DMMR 专用组件
\providecommand{\dmmrReconCard}[3]{% 节点名/标签/样式
  \node[#3, anchor=west, minimum width=16mm, minimum height=10mm, inner sep=0pt] (#1) at (0, 0) {}; % 重建输出外框
  \node[font=\scriptsize, align=center, anchor=north] at ([yshift=-1mm]#1.north) {#2}; % 重建输出标签
  \begin{scope}
    \clip ([xshift=1mm,yshift=1mm]#1.south west) rectangle ([xshift=-1mm,yshift=-3mm]#1.north east); % 限制波形始终在框内
    \begin{scope}[shift={([yshift=2.6mm]#1.south)}, scale=.38] % 重建波形位置
      \pic {eeg signal}; % EEG波形
    \end{scope}
  \end{scope}
  \coordinate (#1-west) at (#1.west); % 左侧连接点
  \coordinate (#1-east) at (#1.east); % 右侧连接点
}

\tikzset{
  % DMMR预处理节点样式
  dmmr prep/.style={
    rounded corners=2pt, % 轻微圆角
    draw=gray!70, % 边框颜色
    fill=gray!8, % 填充颜色
    minimum width=18mm, % 节点宽度
    minimum height=9mm, % 节点高度
    align=center, % 文本居中
    font=\scriptsize % 使用较小字号
  },
  % DMMR注意力节点样式
  dmmr attention/.style={
    rounded corners=2pt, % 轻微圆角
    draw=purple!55!black, % 边框颜色
    fill=purple!12, % 填充颜色
    minimum width=8mm, % 节点宽度
    minimum height=18mm, % 节点高度
    align=center, % 文本居中
    font=\scriptsize % 使用较小字号
  },
  % DMMR编码器节点样式
  dmmr encoder/.style={
    rounded corners=7pt, % 胶囊形圆角
    draw=teal!60!black, % 边框颜色
    fill=cyan!18, % 填充颜色
    minimum width=10mm, % 节点宽度
    minimum height=6mm, % 节点高度
    align=center, % 文本居中
    font=\scriptsize % 使用较小字号
  },
  % DMMR解码器节点样式
  dmmr decoder/.style={
    rounded corners=7pt, % 胶囊形圆角
    draw=green!45!black, % 边框颜色
    fill=green!14, % 填充颜色
    minimum width=12mm, % 节点宽度
    minimum height=8mm, % 节点高度
    align=center, % 文本居中
    font=\scriptsize % 使用较小字号
  },
  % DMMR域判别器节点样式
  dmmr subject discriminator/.style={
    rounded corners=7pt, % 胶囊形圆角
    draw=blue!55!black, % 边框颜色
    fill=blue!18, % 填充颜色
    minimum width=20mm, % 节点宽度
    minimum height=8mm, % 节点高度
    align=center, % 文本居中
    font=\scriptsize % 使用较小字号
  },
  % DMMR阶段背景样式
  dmmr stage/.style={
    rounded corners=4pt, % 背景圆角
    draw=gray!70, % 背景边框
    dashed, % 虚线边框
    fill=green!7, % 默认浅绿背景
    inner sep=3mm % 背景内边距
  },
  % DMMR损失节点样式
  dmmr loss/.style={
    rounded corners=2pt, % 轻微圆角
    draw=orange!75!black, % loss边框
    fill=orange!6, % loss背景
    minimum width=8mm, % 节点宽度
    minimum height=6mm, % 节点高度
    align=center, % 文本居中
    font=\scriptsize % 使用较小字号
  },
  % DMMR总损失样式
  dmmr total loss/.style={
    rounded corners=2pt, % 轻微圆角
    draw=red!70!black, % 总损失边框
    fill=red!10, % 总损失背景
    minimum width=10mm, % 节点宽度
    minimum height=6mm, % 节点高度
    align=center, % 文本居中
    font=\scriptsize % 使用较小字号
  },
  % DMMR混合节点样式
  dmmr mix/.style={
    circle,
    draw=gray!70,
    fill=white,
    minimum size=6mm,
    inner sep=0pt,
    font=\scriptsize
  },
  % DMMR损失相加节点样式
  dmmr sum/.style={
    circle, % 圆形节点
    draw=gray!70, % 边框颜色
    fill=white, % 背景颜色
    minimum size=4mm, % 圆圈大小
    inner sep=0pt, % 无内边距
    font=\scriptsize\bfseries % 加号字号
  },
  % DMMR重建输出样式
  dmmr recon output/.style={
    rounded corners=1.5pt, % 轻微圆角
    draw=gray!50, % 边框颜色
    fill=white, % 背景颜色
    minimum width=16mm, % 节点宽度
    minimum height=11mm, % 节点高度
    inner xsep=1mm, % 横向内边距
    inner ysep=1mm, % 纵向内边距
    align=center, % 文本居中
    font=\scriptsize % 使用较小字号
  },
  % DMMR前向箭头
  dmmr edge/.style={
    -{Latex[length=1.5mm]}, % 箭头样式
    draw=teal!50!black, % 箭头颜色
    line width=.55pt % 箭头线宽
  },
  % DMMR对抗箭头
  dmmr adv edge/.style={
    -{Latex[length=1.5mm]}, % 箭头样式
    draw=purple!65!black, % 对抗分支颜色
    line width=.55pt % 箭头线宽
  },
  % DMMR重建箭头
  dmmr recon edge/.style={
    -{Latex[length=1.5mm]}, % 箭头样式
    draw=green!45!black, % 重建分支颜色
    line width=.55pt % 箭头线宽
  },
  % 小型EEG输出样式
  dmmr mini eeg box/.style={
    rounded corners=1pt, % 轻微圆角
    draw=gray!55, % 边框颜色
    fill=white, % 背景颜色
    minimum width=12mm, % 节点宽度
    minimum height=7mm, % 节点高度
    inner sep=0pt % 无内边距
  },
  % 小型EEG输出组件
  pics/dmmr mini eeg/.style={
    code={ % 定义重建EEG小图
      \node[dmmr mini eeg box] (miniEEGBox) at (0, 0) {}; % 小图外框
      \begin{scope}[scale=.43, transform shape]
        \pic at (0, 0) {eeg signal}; % 复用EEG波形
      \end{scope}
      \coordinate (miniEEG-west) at (miniEEGBox.west); % 左侧连接点
      \coordinate (miniEEG-east) at (miniEEGBox.east); % 右侧连接点
    }
  },
  % DMMR预训练总损失组件
  pics/dmmr pretrain loss/.style={
    code={ % 定义预训练总损失
      \node[dmmr loss] (reconLoss) at (0, .45) {$L_{recon}$}; % 重建损失
      \node[dmmr loss] (advLoss) at (0, -.45) {$L_{adv}$}; % 对抗损失
      \node[
        font=\scriptsize,
        align=center
      ] (preLossText) at (0, -1.20) {$L_{pre}=L_{recon}+L_{adv}$}; % 总预训练损失
      \coordinate (preLoss-recon-west) at (reconLoss.west); % 重建损失入口
      \coordinate (preLoss-adv-west) at (advLoss.west); % 对抗损失入口
    }
  }
}
 % 引入DMMR专用组件

\begin{tikzpicture}[
  node distance=8mm and 11mm
] % DMMR模型图开始
  % 全局间距参数：主要修改这些值即可调整布局
  \def\stepGap{4mm} % 预处理模块横向间距
  \def\preGap{4mm} % 预处理到预训练区域的间距
  \def\mmrGap{4mm} % Encoder到Decoder列、Decoder到重建列的间距
  \def\mixGap{8mm} % Mix到第二阶段Encoder的间距
  \def\mixEncGap{3mm} % 第二阶段Encoder到Decoder列的间距
  \def\mixDecGap{4mm} % 第二阶段多个Decoder的横向间距
  \def\mixReconGap{4mm} % 第二阶段Decoder到recon EEG的间距
  \def\mixBraceGap{1mm} % 第二阶段花括号与Decoder的间距
  \def\mixLabelGap{11mm} % Lrecon标签与花括号的间距
  \def\lossGap{4mm} % 重建输出到loss列的间距
  \def\decoderGap{8mm} % 多个Decoder的纵向间距
  \def\dgdannGap{10mm} % DG-DANN距离重建分支的垂直间距
  \def\dgdannInnerXSep{3mm} % DG-DANN外框横向内边距，pretrain会跟随自适应
  \def\dgdannInnerYSep{4mm} % DG-DANN外框纵向内边距，pretrain会跟随自适应
  \def\preShiftX{0mm} % 整个预训练区域的横向总偏移
  \def\preShiftY{0mm} % 整个预训练区域的纵向总偏移
  \def\fineGap{15mm} % Fine-tuning距离Pre-training的垂直间距
  \def\fineBaseX{-9mm} % Fine-tuning区域的独立横向锚点
  \def\fineBaseY{-40mm} % Fine-tuning区域的独立纵向锚点
  %
  % 输入和预处理流程：从左到右用相对定位排列
  \begin{scope}[shift={(\preShiftX,\preShiftY)}]
  \eegDataCard{rawSample}{Subject i}{dann source data} % 原始EEG样本
  \coordinate (raw-east) at (rawSample.east); % 原始样本连接点
  \node[dmmr prep, right=\stepGap of rawSample] (shuffleNode) {Time-step\\Shuffle}; % 时间步打乱
  \node[dmmr attention, right=\stepGap of shuffleNode] (abpNode) {ABP}; % Attention Based Pooling
  \node[dann source data, minimum width=24mm, minimum height=12mm, right=\stepGap of abpNode] (weightedSample) {}; % 加权样本外框
  \node[font=\scriptsize] at ([yshift=2.4mm]weightedSample.center) {Weighted Sample}; % 加权样本标签
  \begin{scope}[shift={([yshift=-2mm]weightedSample.center)}, scale=.62]
    \pic {eeg signal}; % 加权样本波形
  \end{scope}
  \coordinate (weighted-east) at (weightedSample.east); % 加权样本连接点
  \end{scope}
  %
  % DMMR Pre-training主体：后续模块相对weightedSample定位，因此会跟随preShift整体移动
  \node[dmmr encoder, right=\preGap of weightedSample] (encoderNode) {Encoder}; % 共享编码器
  \node[font=\scriptsize, text=black] (sharedFeatureText) at ([yshift=-4mm]encoderNode.south) {}; % 编码器说明
  \node[dmmr decoder, right=\mmrGap of encoderNode] (decI) {Decoder\\$D_{si}$}; % Subject i解码器
  \node[dmmr decoder, above=\decoderGap of decI] (decOne) {Decoder\\$D_{s1}$}; % Subject 1解码器
  \node[dmmr decoder, below=\decoderGap of decI] (decN) {Decoder\\$D_{sn}$}; % Subject n解码器
  \node[font=\scriptsize] at ($(decOne)!0.5!(decI)$) {$\vdots$}; % Decoder上方省略号
  \node[font=\scriptsize] at ($(decI)!0.5!(decN)$) {$\vdots$}; % Decoder下方省略号
  %
  \node[dmmr mix] (mixNode) at ([xshift=\mixGap]decI.east) {Mix}; % 第一阶段Decoder后接Mix节点
  \node[dmmr encoder, right=\mixEncGap of mixNode] (mixEncoder) {Encoder}; % 第二阶段编码器
  \node[dmmr decoder, right=\mixDecGap of mixEncoder] (mixDecI) {Decoder\\$D_{si}$}; % 第二阶段解码器
  \node[dmmr decoder, above=\decoderGap of mixDecI] (mixDecOne) {Decoder\\$D_{s1}$}; % 二阶段解码器1
  \node[dmmr decoder, below=\decoderGap of mixDecI] (mixDecN) {Decoder\\$D_{sn}$}; % 二阶段解码器n
  \node[font=\scriptsize] at ($(mixDecOne)!0.5!(mixDecI)$) {$\vdots$}; % 二阶段上方省略号
  \node[font=\scriptsize] at ($(mixDecI)!0.5!(mixDecN)$) {$\vdots$}; % 二阶段下方省略号
  \draw[dmmr recon edge] (decOne.east) -- (mixNode.west); % 第一阶段Decoder 1到Mix
  \draw[dmmr recon edge] (decI.east) -- (mixNode.west); % 第一阶段Decoder i到Mix
  \draw[dmmr recon edge] (decN.east) -- (mixNode.west); % 第一阶段Decoder n到Mix
  \draw[dmmr recon edge] (mixNode.east) -- (mixEncoder.west); % Mix到第二阶段Encoder
  \draw[dmmr recon edge] (mixEncoder.east) -- (mixDecOne.west); % 第二阶段Encoder到Decoder 1
  \draw[dmmr recon edge] (mixEncoder.east) -- (mixDecI.west); % 第二阶段Encoder到Decoder i
  \draw[dmmr recon edge] (mixEncoder.east) -- (mixDecN.west); % 第二阶段Encoder到Decoder n
  \begin{scope}[shift={([xshift=\mixReconGap]mixDecOne.east)}]
    \dmmrReconCard{reconOne}{\shortstack{recon $EEG_{1}$}}{dmmr recon output} % 第二阶段recon EEG 1
  \end{scope}
  \begin{scope}[shift={([xshift=\mixReconGap]mixDecI.east)}]
    \dmmrReconCard{reconI}{\shortstack{recon $EEG_{i}$}}{dmmr recon output} % 第二阶段recon EEG i
  \end{scope}
  \begin{scope}[shift={([xshift=\mixReconGap]mixDecN.east)}]
    \dmmrReconCard{reconN}{\shortstack{recon $EEG_{n}$}}{dmmr recon output} % 第二阶段recon EEG n
  \end{scope}
  \draw[dmmr recon edge] (mixDecOne.east) -- (reconOne.west); % 第二阶段Decoder 1到recon EEG 1
  \draw[dmmr recon edge] (mixDecI.east) -- (reconI.west); % 第二阶段Decoder i到recon EEG i
  \draw[dmmr recon edge] (mixDecN.east) -- (reconN.west); % 第二阶段Decoder n到recon EEG n
  \coordinate (reconBraceTop) at ([xshift=\mixBraceGap]reconOne.north east);
  \coordinate (reconBraceBottom) at ([xshift=\mixBraceGap]reconN.south east);
  \draw[
    decorate,
    decoration={brace, mirror, amplitude=3pt},
    draw=gray!70,
    line width=.5pt
  ] (reconBraceBottom) -- (reconBraceTop)
    node[midway, xshift=\mixLabelGap, dmmr loss] (reconLossLabel) {$L_{recon}$}; % 第二阶段recon EEG汇总loss方框
  \node[dmmr sum, right=\lossGap of reconLossLabel] (preSum) {$+$}; % Lrecon和Ladv相加
  \node[dmmr total loss, below=4mm of preSum] (preLoss) {$L_{pre}$}; % 总损失方框
  \coordinate (reconBraceMid) at ($(reconBraceTop)!0.5!(reconBraceBottom)$);
  \draw[dmmr recon edge] ([xshift=1.5mm]reconBraceMid) -- (reconLossLabel.west); % 从花括号右侧精确指向Lrecon的绿色箭头
  %
  % DG-DANN对抗分支：基于Decoder上方锚点整体上移
  \node[dann grl, minimum width=8mm, minimum height=6mm] (grlNode) at ([yshift=\dgdannGap]decOne.north) {GRL}; % 梯度反转层
  \node[dmmr subject discriminator, right=\stepGap of grlNode] (subjectDisc) {Subject\\Discriminator}; % 主体判别器
  \node[dmmr loss, right=\stepGap of subjectDisc] (advLoss) {$L_{adv}$}; % 对抗损失
  %
  % 主流程箭头
  \draw[dmmr edge] (raw-east) -- (shuffleNode.west); % 原始样本到时间步打乱
  \draw[dmmr edge] (shuffleNode.east) -- (abpNode.west); % 打乱样本到ABP
  \draw[dmmr edge] (abpNode.east) -- (weightedSample.west); % ABP到加权样本
  \draw[dmmr edge] (weighted-east) -- (encoderNode.west); % 加权样本到编码器
  %
  % 第一阶段编码器到Decoder分支
  \draw[dmmr recon edge] (encoderNode.east) -- (decOne.west); % 编码特征到Decoder 1
  \draw[dmmr recon edge] (encoderNode.east) -- (decI.west); % 编码特征到Decoder i
  \draw[dmmr recon edge] (encoderNode.east) -- (decN.west); % 编码特征到Decoder n
  \draw[dmmr recon edge] (reconLossLabel.east) -- (preSum.west); % Lrecon进入加法节点
  \draw[dmmr recon edge] (preSum.south) -- (preLoss.north); % 加法节点输出Lpre
  %
  % 对抗分支箭头
  \draw[dmmr adv edge] (encoderNode.north) |- (grlNode.west); % 编码特征进入GRL
  \draw[dmmr adv edge] (grlNode.east) -- (subjectDisc.west); % GRL到主体判别器
  \draw[dmmr adv edge] (subjectDisc.east) -- (advLoss.west); % 主体判别器到对抗损失
  \draw[dmmr adv edge] (advLoss.east) -| (preSum.north); % Ladv进入加法节点
  %
  % 背景分区：pretrainBox 的 fit 节点决定上方子图整体范围；fineInput 的锚点决定下方子图整体位置
  \begin{scope}[on background layer] % 背景层
    \node[
      inner xsep=\dgdannInnerXSep,
      inner ysep=\dgdannInnerYSep,
      fit=(grlNode)(subjectDisc)(advLoss)
    ] (dgdannFit) {}; % DG-DANN外框尺寸锚点，供Pre-training自适应包含
    \node[
      dmmr stage,
      inner xsep=3mm,
      inner ysep=3mm,
      fit=(rawSample)(shuffleNode)(abpNode)(weightedSample)(encoderNode)(sharedFeatureText)(decOne)(decI)(decN)(reconOne)(reconI)(reconN)(mixNode)(mixEncoder)(mixDecOne)(mixDecI)(mixDecN)(reconLossLabel)(preSum)(preLoss)(dgdannFit)
    ] (pretrainBox) {}; % 扩大后的Pre-training区域
    \node[
      dmmr stage,
      draw=gray!60,
      fill=gray!15,
      fill opacity=.90,
      inner xsep=\dgdannInnerXSep,
      inner ysep=\dgdannInnerYSep,
      fit=(grlNode)(subjectDisc)(advLoss)
    ] (dgdannBox) {}; % DG-DANN对抗区域
  \end{scope}
  \node[font=\scriptsize\bfseries, anchor=west] at ([xshift=1mm,yshift=-3mm]pretrainBox.north west) {Pre-training}; % 预训练标题
  \node[font=\fontsize{6pt}{7pt}\selectfont\bfseries, anchor=west] at ([xshift=1mm,yshift=-2mm]dgdannBox.north west) {DG-DANN}; % DG-DANN标题
  %
  % Fine-tuning分支：使用独立锚点，避免跟随pretrainBox一起移动
  \coordinate (fineOrigin) at (\fineBaseX, \fineBaseY); % Fine-tuning独立锚点
  \node[dmmr prep, anchor=west] (fineInput) at (fineOrigin) {EEG\\Sample}; % 微调输入
  \node[dmmr attention, right=\stepGap of fineInput] (fineABP) {ABP}; % 复用注意力层
  \node[dmmr encoder, right=\stepGap of fineABP] (fineEncoder) {Pretrained\\Encoder}; % 复用预训练编码器
  \node[dann head, right=\stepGap of fineEncoder] (fineClassifier) {Classifier\\64$\to$3}; % 微调分类头
  \node[dann output, right=\stepGap of fineClassifier] (finePred) {Predicted\\Label}; % 情绪预测
  \node[font=\scriptsize, text=black, anchor=west] at ([xshift=4mm]finePred.east) {}; % 微调说明
  %
  \draw[dmmr edge] (fineInput.east) -- (fineABP.west); % 微调输入到Attention
  \draw[dmmr edge] (fineABP.east) -- (fineEncoder.west); % Attention到Encoder
  \draw[dmmr edge] (fineEncoder.east) -- (fineClassifier.west); % Encoder到分类层
  \draw[dmmr edge] (fineClassifier.east) -- (finePred.west); % 分类层到情绪预测
  %
  \begin{scope}[on background layer] % 背景层
    \node[
      dmmr stage,
      fill=orange!5,
      inner xsep=3mm,
      inner ysep=3mm,
      fit=(fineInput)(fineABP)(fineEncoder)(fineClassifier)(finePred)
    ] (fineBox) {}; % 微调区域
  \end{scope}
  \node[font=\scriptsize\bfseries, anchor=west] at ([xshift=1mm,yshift=-3mm]fineBox.north west) {Fine-tuning}; % 微调标题
\end{tikzpicture}
 来绘制该模型

% 集中定义 EEG 信号波形模块
\providecommand{\eegDataCard}[3]{ % 节点名/标签/样式
  \node[#3] (#1) at (0, 0) {}; % 数据外框
  \node[font=\scriptsize] at ([yshift=2.4mm]#1.center) {#2}; % 数据标签
  \begin{scope}[shift={([yshift=-2mm]#1.center)}, scale=.62] % EEG波形位置
    \pic {eeg signal}; % EEG波形
  \end{scope}
}
%
\tikzset{
  % EEG白色背景样式
  eeg signal bg/.style={
    draw=gray!25, % 背景边框色
    fill=white, % 背景填充色
    rounded corners=1pt, % 背景轻微圆角
    line width=.25pt % 背景边框宽度
  },
  % EEG波形线样式
  eeg signal line/.style={
    draw=black!65,
    line width=.58pt,
    line cap=round,
    line join=round
  },
  % EEG信号波形
  pics/eeg signal/.style={
    code={ % 定义可复用EEG波形
      \path[eeg signal bg] (-.70, -.26) rectangle (.70, .30); % EEG白色背景
      \draw[eeg signal line]
        (-.62, .00) --
        (-.56, .01) --
        (-.50, -.02) --
        (-.45, .04) --
        (-.40, -.07) --
        (-.35, .05) --
        (-.30, -.02) --
        (-.25, -.15) --
        (-.21, .04) --
        (-.17, .18) --
        (-.13, .06) --
        (-.09, -.06) --
        (-.05, .02) --
        (.00, -.11) --
        (.05, -.04) --
        (.10, -.07) --
        (.16, .03) --
        (.21, -.03) --
        (.26, .08) --
        (.31, -.09) --
        (.36, .02) --
        (.42, -.04) --
        (.48, .03) --
        (.54, -.01) --
        (.58, .00) --
        (.62, .00); % EEG折线
    }
  }
}
 % 引入EEG输入卡片
% 集中定义模型图通用样式和通用组件
\tikzset{
  % 前向传播实线箭头
  dann edge/.style={
    -{Latex[length=1.5mm]}, % 箭头样式
    draw=teal!50!black, % 连线颜色
    line width=.55pt % 连线宽度
  },
  % 域适配/混淆损失虚线箭头
  dann confuse edge/.style={
    -{Latex[length=1.5mm]}, % 箭头样式
    draw=purple!65!black, % 连线颜色
    line width=.55pt, % 连线宽度
    % dashed % 虚线表示域混淆或分布对齐
  },
  % 输入数据节点基础样式
  dann data/.style={
    rounded corners=2pt, % 轻微圆角
    draw=blue!55!black, % 默认边框颜色
    fill=blue!9, % 默认填充颜色
    minimum width=18mm, % 最小宽度
    minimum height=10mm, % 最小高度
    align=center, % 居中换行
    font=\scriptsize % 使用较小字号
  },
  % Source数据节点样式
  dann source data/.style={
    dann data, % 继承数据节点基础样式
    draw=blue!75!black, % 源域边框色
    fill=blue!10 % 源域填充色
  },
  % Target数据节点样式
  dann target data/.style={
    dann data, % 继承数据节点基础样式
    draw=green!50!black, % 目标域边框色
    fill=green!12 % 目标域填充色
  },
  % 分类头节点样式
  dann head/.style={
    rounded corners=2pt, % 轻微圆角
    draw=orange!70!black, % 边框颜色
    fill=orange!13, % 填充颜色
    minimum width=18mm, % 最小宽度
    minimum height=9mm, % 最小高度
    align=center, % 居中换行
    font=\scriptsize % 使用较小字号
  },
  % GRL节点样式
  dann grl/.style={
    rounded corners=2pt, % 轻微圆角
    draw=purple!70!black, % 边框颜色
    fill=purple!12, % 填充颜色
    minimum width=16mm, % 最小宽度
    minimum height=9mm, % 最小高度
    align=center, % 居中换行
    font=\scriptsize % 使用较小加粗字号
  },
  % 域判别器节点样式
  dann domain/.style={
    rounded corners=2pt, % 轻微圆角
    draw=red!65!black, % 边框颜色
    fill=red!10, % 填充颜色
    minimum width=19mm, % 最小宽度
    minimum height=9mm, % 最小高度
    align=center, % 居中换行
    font=\scriptsize % 使用较小字号
  },
  % 输出节点样式
  dann output/.style={
    rounded corners=2pt, % 轻微圆角
    draw=gray!65, % 边框颜色
    fill=gray!8, % 填充颜色
    minimum width=17mm, % 最小宽度
    minimum height=8mm, % 最小高度
    align=center, % 居中换行
    font=\scriptsize % 使用较小字号
  },
  % 简化版共享CFE矩形样式
  dann cfe/.style={
    rounded corners=3pt, % 圆角
    draw=teal!60!black, % 边框颜色
    fill=teal!10, % 填充颜色
    minimum width=24mm, % 最小宽度
    minimum height=16mm, % 最小高度
    align=center, % 居中换行
    font=\scriptsize % 使用较小加粗字号
  },
  % Source EEG 输入组件
  pics/source eeg/.style={
    code={ % 定义Source EEG模块
      \eegDataCard{sourceData}{Source data}{dann source data} % 源域数据卡片
      \coordinate (source-east) at (sourceData.east); % 右侧连接点
    }
  },
  % Target EEG 输入组件
  pics/target eeg/.style={
    code={ % 定义Target EEG模块
      \eegDataCard{targetData}{Target data}{dann target data} % 目标域数据卡片
      \coordinate (target-east) at (targetData.east); % 右侧连接点
    }
  },
  % 简化版共享CFE组件
  pics/shared cfe/.style={
    code={ % 定义共享特征提取器
      \node[dann cfe] (sharedCFE) at (0, 0) {Shared CFE\\310$\to$64}; % 共享CFE矩形
      \node[font=\scriptsize, text=black] at (0, -.92) {Feature Extractor}; % CFE说明
      \coordinate (cfe-source-west) at ([yshift=2.5mm]sharedCFE.west); % 源域输入连接点
      \coordinate (cfe-target-west) at ([yshift=-2.5mm]sharedCFE.west); % 目标域输入连接点
      \coordinate (cfe-emo-east) at ([yshift=2.5mm]sharedCFE.east); % 情绪分支输出点
      \coordinate (cfe-domain-east) at ([yshift=-2.5mm]sharedCFE.east); % 域分支输出点
    }
  },
  % 情绪分类器组件
  pics/emotion classifier/.style={
    code={ % 定义情绪分类器
      \node[dann head] (emoHead) at (0, 0) {Emotion\\Classifier}; % 线性分类层
      \node[dann output] (emoOut) at (2.25, 0) {Predicted\\Label}; % 情感预测
      \draw[dann edge] (emoHead) -- (emoOut); % 分类器到输出
      \coordinate (emo-west) at (emoHead.west); % 左侧连接点
    }
  },
  % GRL组件
  pics/grl/.style={
    code={ % 定义梯度反转层
      \node[dann grl] (grlNode) at (0, 0) {GRL}; % 梯度反转层
      \coordinate (grl-west) at (grlNode.west); % 左侧连接点
      \coordinate (grl-east) at (grlNode.east); % 右侧连接点
    }
  },
  % 域判别器组件
  pics/domain discriminator/.style={
    code={ % 定义域判别器
      \node[dann domain] (domainHead) at (0, 0) {Domain\\Discriminator}; % 域判别器
      \node[dann output] (domainOut) at (2.45, 0) {Source / Target\\Confusion}; % 域混淆输出
      \draw[dann confuse edge] (domainHead) -- (domainOut); % 域判别到混淆
      \coordinate (domain-west) at (domainHead.west); % 左侧连接点
    }
  }
}

% 通用源域/目标域特征图例
\providecommand{\domainFeatureLegend}[2]{ % 锚点节点/下方距离
  \matrix (legendContent) [
    below=#2 of #1,
    column sep=2mm,
    inner sep=0pt,
    nodes={inner sep=0pt, anchor=center},
    ampersand replacement=\&
  ] {
    \node[minimum size=6mm] (legendSourceIcon) {}; \&
    \node[anchor=west, font=\scriptsize, text=gray!95] (legendSourceText) {Source feature}; \&
    \node[minimum width=5mm] {}; \&
    \node[minimum size=6mm] (legendTargetIcon) {}; \&
    \node[anchor=west, font=\scriptsize, text=gray!95] (legendTargetText) {Target feature}; \\
  }; % 图例内容矩阵
  \begin{scope}[shift={(legendSourceIcon.center)}, scale=.30, transform shape, name prefix=legendSource-]
    \pic {source feature cloud}; % 源域特征图例
  \end{scope}
  \begin{scope}[shift={(legendTargetIcon.center)}, scale=.30, transform shape, name prefix=legendTarget-]
    \pic {target feature cloud}; % 目标域特征图例
  \end{scope}
  \begin{scope}[on background layer]
    \node[
      rounded corners=4pt,
      draw=gray!85,
      dashed,
      inner xsep=5mm,
      inner ysep=2mm,
      fit=(legendContent)
    ] (legendBox) {}; % 图例外框
  \end{scope}
}
 % 引入通用节点和箭头样式
% 集中定义 DMMR 专用组件
\providecommand{\dmmrReconCard}[3]{% 节点名/标签/样式
  \node[#3, anchor=west, minimum width=16mm, minimum height=10mm, inner sep=0pt] (#1) at (0, 0) {}; % 重建输出外框
  \node[font=\scriptsize, align=center, anchor=north] at ([yshift=-1mm]#1.north) {#2}; % 重建输出标签
  \begin{scope}
    \clip ([xshift=1mm,yshift=1mm]#1.south west) rectangle ([xshift=-1mm,yshift=-3mm]#1.north east); % 限制波形始终在框内
    \begin{scope}[shift={([yshift=2.6mm]#1.south)}, scale=.38] % 重建波形位置
      \pic {eeg signal}; % EEG波形
    \end{scope}
  \end{scope}
  \coordinate (#1-west) at (#1.west); % 左侧连接点
  \coordinate (#1-east) at (#1.east); % 右侧连接点
}

\tikzset{
  % DMMR预处理节点样式
  dmmr prep/.style={
    rounded corners=2pt, % 轻微圆角
    draw=gray!70, % 边框颜色
    fill=gray!8, % 填充颜色
    minimum width=18mm, % 节点宽度
    minimum height=9mm, % 节点高度
    align=center, % 文本居中
    font=\scriptsize % 使用较小字号
  },
  % DMMR注意力节点样式
  dmmr attention/.style={
    rounded corners=2pt, % 轻微圆角
    draw=purple!55!black, % 边框颜色
    fill=purple!12, % 填充颜色
    minimum width=8mm, % 节点宽度
    minimum height=18mm, % 节点高度
    align=center, % 文本居中
    font=\scriptsize % 使用较小字号
  },
  % DMMR编码器节点样式
  dmmr encoder/.style={
    rounded corners=7pt, % 胶囊形圆角
    draw=teal!60!black, % 边框颜色
    fill=cyan!18, % 填充颜色
    minimum width=10mm, % 节点宽度
    minimum height=6mm, % 节点高度
    align=center, % 文本居中
    font=\scriptsize % 使用较小字号
  },
  % DMMR解码器节点样式
  dmmr decoder/.style={
    rounded corners=7pt, % 胶囊形圆角
    draw=green!45!black, % 边框颜色
    fill=green!14, % 填充颜色
    minimum width=12mm, % 节点宽度
    minimum height=8mm, % 节点高度
    align=center, % 文本居中
    font=\scriptsize % 使用较小字号
  },
  % DMMR域判别器节点样式
  dmmr subject discriminator/.style={
    rounded corners=7pt, % 胶囊形圆角
    draw=blue!55!black, % 边框颜色
    fill=blue!18, % 填充颜色
    minimum width=20mm, % 节点宽度
    minimum height=8mm, % 节点高度
    align=center, % 文本居中
    font=\scriptsize % 使用较小字号
  },
  % DMMR阶段背景样式
  dmmr stage/.style={
    rounded corners=4pt, % 背景圆角
    draw=gray!70, % 背景边框
    dashed, % 虚线边框
    fill=green!7, % 默认浅绿背景
    inner sep=3mm % 背景内边距
  },
  % DMMR损失节点样式
  dmmr loss/.style={
    rounded corners=2pt, % 轻微圆角
    draw=orange!75!black, % loss边框
    fill=orange!6, % loss背景
    minimum width=8mm, % 节点宽度
    minimum height=6mm, % 节点高度
    align=center, % 文本居中
    font=\scriptsize % 使用较小字号
  },
  % DMMR总损失样式
  dmmr total loss/.style={
    rounded corners=2pt, % 轻微圆角
    draw=red!70!black, % 总损失边框
    fill=red!10, % 总损失背景
    minimum width=10mm, % 节点宽度
    minimum height=6mm, % 节点高度
    align=center, % 文本居中
    font=\scriptsize % 使用较小字号
  },
  % DMMR混合节点样式
  dmmr mix/.style={
    circle,
    draw=gray!70,
    fill=white,
    minimum size=6mm,
    inner sep=0pt,
    font=\scriptsize
  },
  % DMMR损失相加节点样式
  dmmr sum/.style={
    circle, % 圆形节点
    draw=gray!70, % 边框颜色
    fill=white, % 背景颜色
    minimum size=4mm, % 圆圈大小
    inner sep=0pt, % 无内边距
    font=\scriptsize\bfseries % 加号字号
  },
  % DMMR重建输出样式
  dmmr recon output/.style={
    rounded corners=1.5pt, % 轻微圆角
    draw=gray!50, % 边框颜色
    fill=white, % 背景颜色
    minimum width=16mm, % 节点宽度
    minimum height=11mm, % 节点高度
    inner xsep=1mm, % 横向内边距
    inner ysep=1mm, % 纵向内边距
    align=center, % 文本居中
    font=\scriptsize % 使用较小字号
  },
  % DMMR前向箭头
  dmmr edge/.style={
    -{Latex[length=1.5mm]}, % 箭头样式
    draw=teal!50!black, % 箭头颜色
    line width=.55pt % 箭头线宽
  },
  % DMMR对抗箭头
  dmmr adv edge/.style={
    -{Latex[length=1.5mm]}, % 箭头样式
    draw=purple!65!black, % 对抗分支颜色
    line width=.55pt % 箭头线宽
  },
  % DMMR重建箭头
  dmmr recon edge/.style={
    -{Latex[length=1.5mm]}, % 箭头样式
    draw=green!45!black, % 重建分支颜色
    line width=.55pt % 箭头线宽
  },
  % 小型EEG输出样式
  dmmr mini eeg box/.style={
    rounded corners=1pt, % 轻微圆角
    draw=gray!55, % 边框颜色
    fill=white, % 背景颜色
    minimum width=12mm, % 节点宽度
    minimum height=7mm, % 节点高度
    inner sep=0pt % 无内边距
  },
  % 小型EEG输出组件
  pics/dmmr mini eeg/.style={
    code={ % 定义重建EEG小图
      \node[dmmr mini eeg box] (miniEEGBox) at (0, 0) {}; % 小图外框
      \begin{scope}[scale=.43, transform shape]
        \pic at (0, 0) {eeg signal}; % 复用EEG波形
      \end{scope}
      \coordinate (miniEEG-west) at (miniEEGBox.west); % 左侧连接点
      \coordinate (miniEEG-east) at (miniEEGBox.east); % 右侧连接点
    }
  },
  % DMMR预训练总损失组件
  pics/dmmr pretrain loss/.style={
    code={ % 定义预训练总损失
      \node[dmmr loss] (reconLoss) at (0, .45) {$L_{recon}$}; % 重建损失
      \node[dmmr loss] (advLoss) at (0, -.45) {$L_{adv}$}; % 对抗损失
      \node[
        font=\scriptsize,
        align=center
      ] (preLossText) at (0, -1.20) {$L_{pre}=L_{recon}+L_{adv}$}; % 总预训练损失
      \coordinate (preLoss-recon-west) at (reconLoss.west); % 重建损失入口
      \coordinate (preLoss-adv-west) at (advLoss.west); % 对抗损失入口
    }
  }
}
 % 引入DMMR专用组件

\begin{tikzpicture}[
  node distance=8mm and 11mm
] % DMMR模型图开始
  % 全局间距参数：主要修改这些值即可调整布局
  \def\stepGap{4mm} % 预处理模块横向间距
  \def\preGap{4mm} % 预处理到预训练区域的间距
  \def\mmrGap{4mm} % Encoder到Decoder列、Decoder到重建列的间距
  \def\mixGap{8mm} % Mix到第二阶段Encoder的间距
  \def\mixEncGap{3mm} % 第二阶段Encoder到Decoder列的间距
  \def\mixDecGap{4mm} % 第二阶段多个Decoder的横向间距
  \def\mixReconGap{4mm} % 第二阶段Decoder到recon EEG的间距
  \def\mixBraceGap{1mm} % 第二阶段花括号与Decoder的间距
  \def\mixLabelGap{11mm} % Lrecon标签与花括号的间距
  \def\lossGap{4mm} % 重建输出到loss列的间距
  \def\decoderGap{8mm} % 多个Decoder的纵向间距
  \def\dgdannGap{10mm} % DG-DANN距离重建分支的垂直间距
  \def\dgdannInnerXSep{3mm} % DG-DANN外框横向内边距，pretrain会跟随自适应
  \def\dgdannInnerYSep{4mm} % DG-DANN外框纵向内边距，pretrain会跟随自适应
  \def\preShiftX{0mm} % 整个预训练区域的横向总偏移
  \def\preShiftY{0mm} % 整个预训练区域的纵向总偏移
  \def\fineGap{15mm} % Fine-tuning距离Pre-training的垂直间距
  \def\fineBaseX{-9mm} % Fine-tuning区域的独立横向锚点
  \def\fineBaseY{-40mm} % Fine-tuning区域的独立纵向锚点
  %
  % 输入和预处理流程：从左到右用相对定位排列
  \begin{scope}[shift={(\preShiftX,\preShiftY)}]
  \eegDataCard{rawSample}{Subject i}{dann source data} % 原始EEG样本
  \coordinate (raw-east) at (rawSample.east); % 原始样本连接点
  \node[dmmr prep, right=\stepGap of rawSample] (shuffleNode) {Time-step\\Shuffle}; % 时间步打乱
  \node[dmmr attention, right=\stepGap of shuffleNode] (abpNode) {ABP}; % Attention Based Pooling
  \node[dann source data, minimum width=24mm, minimum height=12mm, right=\stepGap of abpNode] (weightedSample) {}; % 加权样本外框
  \node[font=\scriptsize] at ([yshift=2.4mm]weightedSample.center) {Weighted Sample}; % 加权样本标签
  \begin{scope}[shift={([yshift=-2mm]weightedSample.center)}, scale=.62]
    \pic {eeg signal}; % 加权样本波形
  \end{scope}
  \coordinate (weighted-east) at (weightedSample.east); % 加权样本连接点
  \end{scope}
  %
  % DMMR Pre-training主体：后续模块相对weightedSample定位，因此会跟随preShift整体移动
  \node[dmmr encoder, right=\preGap of weightedSample] (encoderNode) {Encoder}; % 共享编码器
  \node[font=\scriptsize, text=black] (sharedFeatureText) at ([yshift=-4mm]encoderNode.south) {}; % 编码器说明
  \node[dmmr decoder, right=\mmrGap of encoderNode] (decI) {Decoder\\$D_{si}$}; % Subject i解码器
  \node[dmmr decoder, above=\decoderGap of decI] (decOne) {Decoder\\$D_{s1}$}; % Subject 1解码器
  \node[dmmr decoder, below=\decoderGap of decI] (decN) {Decoder\\$D_{sn}$}; % Subject n解码器
  \node[font=\scriptsize] at ($(decOne)!0.5!(decI)$) {$\vdots$}; % Decoder上方省略号
  \node[font=\scriptsize] at ($(decI)!0.5!(decN)$) {$\vdots$}; % Decoder下方省略号
  %
  \node[dmmr mix] (mixNode) at ([xshift=\mixGap]decI.east) {Mix}; % 第一阶段Decoder后接Mix节点
  \node[dmmr encoder, right=\mixEncGap of mixNode] (mixEncoder) {Encoder}; % 第二阶段编码器
  \node[dmmr decoder, right=\mixDecGap of mixEncoder] (mixDecI) {Decoder\\$D_{si}$}; % 第二阶段解码器
  \node[dmmr decoder, above=\decoderGap of mixDecI] (mixDecOne) {Decoder\\$D_{s1}$}; % 二阶段解码器1
  \node[dmmr decoder, below=\decoderGap of mixDecI] (mixDecN) {Decoder\\$D_{sn}$}; % 二阶段解码器n
  \node[font=\scriptsize] at ($(mixDecOne)!0.5!(mixDecI)$) {$\vdots$}; % 二阶段上方省略号
  \node[font=\scriptsize] at ($(mixDecI)!0.5!(mixDecN)$) {$\vdots$}; % 二阶段下方省略号
  \draw[dmmr recon edge] (decOne.east) -- (mixNode.west); % 第一阶段Decoder 1到Mix
  \draw[dmmr recon edge] (decI.east) -- (mixNode.west); % 第一阶段Decoder i到Mix
  \draw[dmmr recon edge] (decN.east) -- (mixNode.west); % 第一阶段Decoder n到Mix
  \draw[dmmr recon edge] (mixNode.east) -- (mixEncoder.west); % Mix到第二阶段Encoder
  \draw[dmmr recon edge] (mixEncoder.east) -- (mixDecOne.west); % 第二阶段Encoder到Decoder 1
  \draw[dmmr recon edge] (mixEncoder.east) -- (mixDecI.west); % 第二阶段Encoder到Decoder i
  \draw[dmmr recon edge] (mixEncoder.east) -- (mixDecN.west); % 第二阶段Encoder到Decoder n
  \begin{scope}[shift={([xshift=\mixReconGap]mixDecOne.east)}]
    \dmmrReconCard{reconOne}{\shortstack{recon $EEG_{1}$}}{dmmr recon output} % 第二阶段recon EEG 1
  \end{scope}
  \begin{scope}[shift={([xshift=\mixReconGap]mixDecI.east)}]
    \dmmrReconCard{reconI}{\shortstack{recon $EEG_{i}$}}{dmmr recon output} % 第二阶段recon EEG i
  \end{scope}
  \begin{scope}[shift={([xshift=\mixReconGap]mixDecN.east)}]
    \dmmrReconCard{reconN}{\shortstack{recon $EEG_{n}$}}{dmmr recon output} % 第二阶段recon EEG n
  \end{scope}
  \draw[dmmr recon edge] (mixDecOne.east) -- (reconOne.west); % 第二阶段Decoder 1到recon EEG 1
  \draw[dmmr recon edge] (mixDecI.east) -- (reconI.west); % 第二阶段Decoder i到recon EEG i
  \draw[dmmr recon edge] (mixDecN.east) -- (reconN.west); % 第二阶段Decoder n到recon EEG n
  \coordinate (reconBraceTop) at ([xshift=\mixBraceGap]reconOne.north east);
  \coordinate (reconBraceBottom) at ([xshift=\mixBraceGap]reconN.south east);
  \draw[
    decorate,
    decoration={brace, mirror, amplitude=3pt},
    draw=gray!70,
    line width=.5pt
  ] (reconBraceBottom) -- (reconBraceTop)
    node[midway, xshift=\mixLabelGap, dmmr loss] (reconLossLabel) {$L_{recon}$}; % 第二阶段recon EEG汇总loss方框
  \node[dmmr sum, right=\lossGap of reconLossLabel] (preSum) {$+$}; % Lrecon和Ladv相加
  \node[dmmr total loss, below=4mm of preSum] (preLoss) {$L_{pre}$}; % 总损失方框
  \coordinate (reconBraceMid) at ($(reconBraceTop)!0.5!(reconBraceBottom)$);
  \draw[dmmr recon edge] ([xshift=1.5mm]reconBraceMid) -- (reconLossLabel.west); % 从花括号右侧精确指向Lrecon的绿色箭头
  %
  % DG-DANN对抗分支：基于Decoder上方锚点整体上移
  \node[dann grl, minimum width=8mm, minimum height=6mm] (grlNode) at ([yshift=\dgdannGap]decOne.north) {GRL}; % 梯度反转层
  \node[dmmr subject discriminator, right=\stepGap of grlNode] (subjectDisc) {Subject\\Discriminator}; % 主体判别器
  \node[dmmr loss, right=\stepGap of subjectDisc] (advLoss) {$L_{adv}$}; % 对抗损失
  %
  % 主流程箭头
  \draw[dmmr edge] (raw-east) -- (shuffleNode.west); % 原始样本到时间步打乱
  \draw[dmmr edge] (shuffleNode.east) -- (abpNode.west); % 打乱样本到ABP
  \draw[dmmr edge] (abpNode.east) -- (weightedSample.west); % ABP到加权样本
  \draw[dmmr edge] (weighted-east) -- (encoderNode.west); % 加权样本到编码器
  %
  % 第一阶段编码器到Decoder分支
  \draw[dmmr recon edge] (encoderNode.east) -- (decOne.west); % 编码特征到Decoder 1
  \draw[dmmr recon edge] (encoderNode.east) -- (decI.west); % 编码特征到Decoder i
  \draw[dmmr recon edge] (encoderNode.east) -- (decN.west); % 编码特征到Decoder n
  \draw[dmmr recon edge] (reconLossLabel.east) -- (preSum.west); % Lrecon进入加法节点
  \draw[dmmr recon edge] (preSum.south) -- (preLoss.north); % 加法节点输出Lpre
  %
  % 对抗分支箭头
  \draw[dmmr adv edge] (encoderNode.north) |- (grlNode.west); % 编码特征进入GRL
  \draw[dmmr adv edge] (grlNode.east) -- (subjectDisc.west); % GRL到主体判别器
  \draw[dmmr adv edge] (subjectDisc.east) -- (advLoss.west); % 主体判别器到对抗损失
  \draw[dmmr adv edge] (advLoss.east) -| (preSum.north); % Ladv进入加法节点
  %
  % 背景分区：pretrainBox 的 fit 节点决定上方子图整体范围；fineInput 的锚点决定下方子图整体位置
  \begin{scope}[on background layer] % 背景层
    \node[
      inner xsep=\dgdannInnerXSep,
      inner ysep=\dgdannInnerYSep,
      fit=(grlNode)(subjectDisc)(advLoss)
    ] (dgdannFit) {}; % DG-DANN外框尺寸锚点，供Pre-training自适应包含
    \node[
      dmmr stage,
      inner xsep=3mm,
      inner ysep=3mm,
      fit=(rawSample)(shuffleNode)(abpNode)(weightedSample)(encoderNode)(sharedFeatureText)(decOne)(decI)(decN)(reconOne)(reconI)(reconN)(mixNode)(mixEncoder)(mixDecOne)(mixDecI)(mixDecN)(reconLossLabel)(preSum)(preLoss)(dgdannFit)
    ] (pretrainBox) {}; % 扩大后的Pre-training区域
    \node[
      dmmr stage,
      draw=gray!60,
      fill=gray!15,
      fill opacity=.90,
      inner xsep=\dgdannInnerXSep,
      inner ysep=\dgdannInnerYSep,
      fit=(grlNode)(subjectDisc)(advLoss)
    ] (dgdannBox) {}; % DG-DANN对抗区域
  \end{scope}
  \node[font=\scriptsize\bfseries, anchor=west] at ([xshift=1mm,yshift=-3mm]pretrainBox.north west) {Pre-training}; % 预训练标题
  \node[font=\fontsize{6pt}{7pt}\selectfont\bfseries, anchor=west] at ([xshift=1mm,yshift=-2mm]dgdannBox.north west) {DG-DANN}; % DG-DANN标题
  %
  % Fine-tuning分支：使用独立锚点，避免跟随pretrainBox一起移动
  \coordinate (fineOrigin) at (\fineBaseX, \fineBaseY); % Fine-tuning独立锚点
  \node[dmmr prep, anchor=west] (fineInput) at (fineOrigin) {EEG\\Sample}; % 微调输入
  \node[dmmr attention, right=\stepGap of fineInput] (fineABP) {ABP}; % 复用注意力层
  \node[dmmr encoder, right=\stepGap of fineABP] (fineEncoder) {Pretrained\\Encoder}; % 复用预训练编码器
  \node[dann head, right=\stepGap of fineEncoder] (fineClassifier) {Classifier\\64$\to$3}; % 微调分类头
  \node[dann output, right=\stepGap of fineClassifier] (finePred) {Predicted\\Label}; % 情绪预测
  \node[font=\scriptsize, text=black, anchor=west] at ([xshift=4mm]finePred.east) {}; % 微调说明
  %
  \draw[dmmr edge] (fineInput.east) -- (fineABP.west); % 微调输入到Attention
  \draw[dmmr edge] (fineABP.east) -- (fineEncoder.west); % Attention到Encoder
  \draw[dmmr edge] (fineEncoder.east) -- (fineClassifier.west); % Encoder到分类层
  \draw[dmmr edge] (fineClassifier.east) -- (finePred.west); % 分类层到情绪预测
  %
  \begin{scope}[on background layer] % 背景层
    \node[
      dmmr stage,
      fill=orange!5,
      inner xsep=3mm,
      inner ysep=3mm,
      fit=(fineInput)(fineABP)(fineEncoder)(fineClassifier)(finePred)
    ] (fineBox) {}; % 微调区域
  \end{scope}
  \node[font=\scriptsize\bfseries, anchor=west] at ([xshift=1mm,yshift=-3mm]fineBox.north west) {Fine-tuning}; % 微调标题
\end{tikzpicture}
